% =========================================================
% Chapter: Electrostatics
% POSN 1 Physics Preparation Notes
% Language: Thai
% Compiler: XeLaTeX
% =========================================================

\chapter{ไฟฟ้าสถิต}

\begin{center}
    {\Large \textbf{Electrostatics}}\\[4pt]
    {\normalsize เอกสารเตรียมสอบคัดเลือก สอวน. ฟิสิกส์ ค่าย 1}\\[6pt]
    {\normalsize จัดทำโดย \textbf{Tames Thanakrit Damduan}}\\[2pt]
    {\footnotesize \textcopyright\ 2026 Tames Thanakrit Damduan. All rights reserved.}\\[8pt]
\end{center}

\section*{เป้าหมายของบทเรียน}

หลังจากเรียนบทนี้ นักเรียนควรทำสิ่งต่อไปนี้ได้

\begin{enumerate}
    \item เข้าใจประจุไฟฟ้า การอนุรักษ์ประจุ และการเหนี่ยวนำไฟฟ้าเบื้องต้น
    \item ใช้กฎของคูลอมบ์ในการหาแรงระหว่างประจุจุดได้
    \item ใช้หลักการซ้อนทับเพื่อหาแรงลัพธ์และสนามไฟฟ้าลัพธ์ได้
    \item เข้าใจสนามไฟฟ้าและเส้นสนามไฟฟ้าในเชิงเวกเตอร์
    \item เข้าใจฟลักซ์ไฟฟ้า เลือกพื้นผิวเกาส์ที่เหมาะสม และใช้กฎของเกาส์กับระบบที่มีสมมาตรได้
    \item หางานของแรงไฟฟ้า พลังงานศักย์ไฟฟ้า และศักย์ไฟฟ้าได้
    \item ใช้ความสัมพันธ์ระหว่างสนามไฟฟ้าและศักย์ไฟฟ้าในกรณีพื้นฐานได้
    \item อธิบายและพิสูจน์ความจุของตัวเก็บประจุแผ่นขนาน การต่ออนุกรม การต่อขนาน พลังงาน และแรงระหว่างแผ่นได้
    \item แก้โจทย์ไฟฟ้าสถิตระดับ สอวน. ที่ผสมเวกเตอร์ สมดุล งาน พลังงาน และตัวเก็บประจุได้
\end{enumerate}

\begin{tcolorbox}[title=แนวคิดสำคัญของบทนี้]
ไฟฟ้าสถิตไม่ได้ยากเพราะสูตรเยอะ แต่ยากเพราะต้องแยกให้ชัดว่าอะไรเป็นสเกลาร์และอะไรเป็นเวกเตอร์ แรงไฟฟ้าและสนามไฟฟ้าเป็นเวกเตอร์ ต้องรวมแบบเวกเตอร์ ส่วนศักย์ไฟฟ้าและพลังงานศักย์ไฟฟ้าเป็นสเกลาร์ จึงรวมแบบพีชคณิตได้โดยตรง
\end{tcolorbox}

\newpage{}

% =========================================================
\section{ประจุไฟฟ้าและกฎพื้นฐาน}

\subsection{ประจุไฟฟ้า}

ประจุไฟฟ้าเป็นสมบัติพื้นฐานของสสาร มีสองชนิดคือ

\[
\text{ประจุบวก}
\]

และ

\[
\text{ประจุลบ}
\]

ประจุชนิดเดียวกันผลักกัน ส่วนประจุต่างชนิดกันดึงดูดกัน

\begin{tcolorbox}[title=นิยาม]
ประจุไฟฟ้าเป็นปริมาณสเกลาร์ที่มีเครื่องหมาย หน่วย SI คือคูลอมบ์ เขียนเป็น
\[
\si{C}
\]
\end{tcolorbox}

ประจุของอิเล็กตรอนมีค่า

\[
-e=-1.60\times10^{-19}\,\si{C}
\]

และประจุของโปรตอนมีค่า

\[
+e=+1.60\times10^{-19}\,\si{C}
\]

\subsection{การควอนไทซ์ของประจุ}

ประจุรวมของวัตถุหนึ่ง ๆ เป็นจำนวนเต็มเท่าของประจุมูลฐาน

\[
q=ne
\]

โดยที่ $n$ เป็นจำนวนเต็มบวก ลบ หรือศูนย์

\subsection{การอนุรักษ์ประจุ}

ในระบบปิด ประจุไฟฟ้ารวมไม่เปลี่ยน

\[
q_{\text{total, before}}=q_{\text{total, after}}
\]

ประจุสามารถถ่ายโอนจากวัตถุหนึ่งไปอีกวัตถุหนึ่งได้ แต่ไม่หายไปเองและไม่เกิดขึ้นเองแบบโดด ๆ

\begin{tcolorbox}[title=ข้อควรจำ]
เวลาโจทย์บอกว่าเอาวัตถุนำไฟฟ้าสองอันมาแตะกัน ประจุรวมต้องคงเดิม แต่ประจุบนแต่ละวัตถุอาจเปลี่ยนจนศักย์ไฟฟ้าของตัวนำเท่ากัน
\end{tcolorbox}

% =========================================================
\section{กฎของคูลอมบ์}

\subsection{แรงระหว่างประจุจุดสองประจุ}

ถ้ามีประจุจุดสองประจุ $q_1$ และ $q_2$ อยู่ห่างกันเป็นระยะ $r$ ขนาดของแรงไฟฟ้าระหว่างกันคือ

\[
F=k_e\frac{|q_1q_2|}{r^2}
\]

โดย

\[
k_e=\frac{1}{4\pi\varepsilon_0}
\]

และ

\[
\varepsilon_0=8.85\times10^{-12}\,\si{C^2/(N.m^2)}
\]

\begin{tcolorbox}[title=กฎของคูลอมบ์]
\[
F=k_e\frac{|q_1q_2|}{r^2}
\]

แรงอยู่ตามแนวเส้นตรงที่เชื่อมประจุทั้งสอง

ถ้า $q_1q_2>0$ แรงเป็นแรงผลัก

ถ้า $q_1q_2<0$ แรงเป็นแรงดึงดูด
\end{tcolorbox}

\subsection{รูปเวกเตอร์ของแรงไฟฟ้า}

ถ้าต้องการหาแรงที่ประจุ $q_2$ กระทำต่อ $q_1$ ให้ใช้แนวคิดว่าแรงชี้ไปตามแนวรัศมีระหว่างประจุ

\[
\vec{F}_{12}=k_e\frac{q_1q_2}{r^2}\hat{r}_{12}
\]

โดย $\hat{r}_{12}$ คือเวกเตอร์หน่วยที่เลือกตามแนวจากประจุหนึ่งไปอีกประจุหนึ่งตามนิยามที่ใช้ในโจทย์

ในระดับ สอวน. ค่าย 1 วิธีที่ปลอดภัยคือ

\begin{enumerate}
    \item หาขนาดแรงด้วย $F=k_e|q_1q_2|/r^2$
    \item ตัดสินทิศจากชนิดของประจุ
    \item แตกแรงเป็นองค์ประกอบตามแกนถ้าจำเป็น
\end{enumerate}

\subsection{หลักการซ้อนทับของแรง}

ถ้าประจุหนึ่งถูกแรงจากหลายประจุพร้อมกัน แรงลัพธ์คือผลรวมเวกเตอร์ของแรงทั้งหมด

\[
\vec{F}_{\text{net}}=\vec{F}_1+\vec{F}_2+\vec{F}_3+\cdots
\]

\begin{tcolorbox}[title=ข้อควรจำสำหรับโจทย์หลายประจุ]
แรงไฟฟ้าเป็นเวกเตอร์ จึงต้องรวมแบบเวกเตอร์ ถ้าแรงไม่อยู่แนวเดียวกัน ให้แตกเป็นแกน $x$ และ $y$ ก่อนเสมอ
\end{tcolorbox}

\begin{examplebox}{ตัวอย่าง: แรงไฟฟ้าบนแกนเดียวกัน}
ประจุ $+q$ อยู่ที่ $x=0$ และประจุ $+4q$ อยู่ที่ $x=3a$ ต้องการหาตำแหน่งบนแกน $x$ ที่วางประจุทดสอบแล้วแรงลัพธ์เป็นศูนย์

เพราะประจุทั้งสองเป็นบวก ตำแหน่งที่แรงหักล้างกันต้องอยู่ระหว่างสองประจุ

ให้จุดนั้นอยู่ห่างจาก $+q$ เป็นระยะ $x$ ดังนั้นห่างจาก $+4q$ เป็นระยะ $3a-x$

ขนาดแรงจากสองข้างต้องเท่ากัน

\[
k_e\frac{qQ}{x^2}=k_e\frac{4qQ}{(3a-x)^2}
\]

ตัด $k_e qQ$

\[
\frac{1}{x^2}=\frac{4}{(3a-x)^2}
\]

ถอดราก

\[
\frac{1}{x}=\frac{2}{3a-x}
\]

\[
3a-x=2x
\]

\[
x=a
\]

ดังนั้นจุดสมดุลอยู่ห่างจากประจุ $+q$ เป็นระยะ

\[
\boxed{a}
\]
\end{examplebox}

\begin{exercisebox}{แบบฝึกหัด 1: กฎของคูลอมบ์}
\begin{enumerate}
    \item ประจุ $+2q$ และ $-q$ อยู่ห่างกัน $r$ จงหาขนาดแรงไฟฟ้าระหว่างกัน และบอกว่าเป็นแรงดึงดูดหรือผลัก
    \item ประจุ $+q$ อยู่ที่ $x=0$ และประจุ $+9q$ อยู่ที่ $x=4a$ จงหาตำแหน่งระหว่างสองประจุที่สนามไฟฟ้ารวมเป็นศูนย์
    \item ประจุ $+q$ สามตัวอยู่ที่มุมของสามเหลี่ยมด้านเท่าด้านยาว $a$ จงหาขนาดแรงลัพธ์บนประจุที่มุมหนึ่ง
    \item ประจุ $+q$ และ $-q$ อยู่ที่ $x=-a$ และ $x=+a$ ตามลำดับ จงพิจารณาทิศของแรงบนประจุทดสอบบวกที่จุด $(0,b)$
\end{enumerate}
\end{exercisebox}

% =========================================================
\section{สนามไฟฟ้า}

\subsection{นิยามของสนามไฟฟ้า}

สนามไฟฟ้าบอกว่า ถ้าเอาประจุทดสอบบวกขนาดเล็กมากไปวางที่จุดหนึ่ง จะถูกแรงต่อหนึ่งหน่วยประจุเท่าไร

\[
\vec{E}=\frac{\vec{F}}{q_0}
\]

ดังนั้น

\[
\vec{F}=q_0\vec{E}
\]

ถ้า $q_0>0$ แรงมีทิศเดียวกับสนามไฟฟ้า

ถ้า $q_0<0$ แรงมีทิศตรงข้ามกับสนามไฟฟ้า

\begin{tcolorbox}[title=นิยามสนามไฟฟ้า]
\[
\vec{E}=\frac{\vec{F}}{q_0}
\]

หน่วยของสนามไฟฟ้าคือ

\[
\si{N/C}
\]

หรือเท่ากับ

\[
\si{V/m}
\]
\end{tcolorbox}

\subsection{สนามไฟฟ้าจากประจุจุด}

ประจุจุด $q$ ทำให้เกิดสนามไฟฟ้าที่จุดห่างออกไป $r$ มีขนาด

\[
E=k_e\frac{|q|}{r^2}
\]

ทิศทางของสนามคือ

\[
\text{ชี้ออกจากประจุบวก}
\]

และ

\[
\text{ชี้เข้าหาประจุลบ}
\]

\subsection{หลักการซ้อนทับของสนามไฟฟ้า}

สนามไฟฟ้าลัพธ์จากหลายประจุคือผลรวมเวกเตอร์ของสนามจากแต่ละประจุ

\[
\vec{E}_{\text{net}}=\vec{E}_1+\vec{E}_2+\vec{E}_3+\cdots
\]

\begin{tcolorbox}[title=ข้อควรจำ]
สนามไฟฟ้าเป็นเวกเตอร์ ดังนั้นต้องดูทิศและรวมแบบเวกเตอร์เหมือนแรง แต่สนามไฟฟ้าไม่ต้องคูณประจุทดสอบก่อนรวม
\end{tcolorbox}

\subsection{สนามไฟฟ้าจากวงแหวนประจุบนแกนสมมาตร}

ข้อสอบ สอวน. มักใช้สมมาตรของวงแหวนประจุ เพราะองค์ประกอบแนวขวางหักล้างกันหมด เหลือเฉพาะองค์ประกอบตามแกน

ให้วงแหวนรัศมี $R$ มีประจุรวม $Q$ กระจายสม่ำเสมอ จุด $P$ อยู่บนแกนของวงแหวน ห่างจากศูนย์กลาง $x$

ระยะจากชิ้นประจุเล็ก ๆ บนวงแหวนถึงจุด $P$ เท่ากันหมดคือ

\[
r=\sqrt{R^2+x^2}
\]

สนามจากชิ้นประจุเล็ก ๆ มีองค์ประกอบตามแกนเท่านั้นที่รวมกันได้

\[
dE_x=dE\cos\theta
\]

โดย

\[
\cos\theta=\frac{x}{\sqrt{R^2+x^2}}
\]

และ

\[
dE=k_e\frac{dq}{R^2+x^2}
\]

จึงได้

\[
dE_x=k_e\frac{x\,dq}{(R^2+x^2)^{3/2}}
\]

รวมทั้งวงแหวน

\[
E_x=k_e\frac{x}{(R^2+x^2)^{3/2}}\int dq
\]

\[
E_x=k_e\frac{Qx}{(R^2+x^2)^{3/2}}
\]

\begin{tcolorbox}[title=สนามไฟฟ้าบนแกนของวงแหวนประจุ]
\[
E_x=k_e\frac{Qx}{(R^2+x^2)^{3/2}}
\]

ถ้า $Q>0$ และ $x>0$ สนามชี้ไปตามแกนบวก
\end{tcolorbox}

\subsection{เส้นสนามไฟฟ้า}

เส้นสนามไฟฟ้าใช้แสดงทิศของสนามไฟฟ้า

\begin{enumerate}
    \item เส้นสนามออกจากประจุบวกและเข้าสู่ประจุลบ
    \item จุดที่เส้นสนามถี่กว่า มีสนามไฟฟ้าแรงกว่า
    \item เส้นสนามไม่ตัดกัน
    \item ทิศของสนามที่จุดหนึ่งคือทิศสัมผัสกับเส้นสนามที่จุดนั้น
\end{enumerate}

\begin{examplebox}{ตัวอย่าง: สนามไฟฟ้าจากประจุสองตัวสมมาตร}
ประจุ $+q$ อยู่ที่ $(0,a)$ และ $+q$ อยู่ที่ $(0,-a)$ ต้องการหาทิศของสนามไฟฟ้าที่จุด $(x,0)$ เมื่อ $x>0$

สนามจากประจุบนและล่างมีขนาดเท่ากัน เพราะระยะถึงจุด $(x,0)$ เท่ากัน

องค์ประกอบตามแกน $y$ ของสนามทั้งสองหักล้างกัน

\[
E_y=0
\]

องค์ประกอบตามแกน $x$ ชี้ไปทางขวาทั้งคู่ จึงรวมกัน

\[
E_x>0
\]

ดังนั้นสนามลัพธ์ชี้ไปทางแกน $+x$
\end{examplebox}

\begin{exercisebox}{แบบฝึกหัด 2: สนามไฟฟ้า}
\begin{enumerate}
    \item ประจุ $+Q$ อยู่ที่จุดกำเนิด จงหาขนาดและทิศของสนามไฟฟ้าที่จุดห่างออกไป $r$ บนแกน $+x$
    \item ประจุ $-Q$ อยู่ที่จุดกำเนิด จงหาขนาดและทิศของสนามไฟฟ้าที่จุดห่างออกไป $r$ บนแกน $+x$
    \item ประจุ $+q$ สองตัวอยู่ที่ $(0,a)$ และ $(0,-a)$ จงหาทิศของสนามไฟฟ้าที่จุด $(x,0)$
    \item วงแหวนรัศมี $R$ มีประจุรวม $Q$ จงเขียนสูตรสนามไฟฟ้าบนแกนของวงแหวนที่ระยะ $x$ จากศูนย์กลาง
\end{enumerate}
\end{exercisebox}

% =========================================================
\section{ฟลักซ์ไฟฟ้าและกฎของเกาส์}

\subsection{แนวคิดของฟลักซ์ไฟฟ้า}

สนามไฟฟ้าบอกขนาดและทิศของแรงต่อหนึ่งหน่วยประจุที่แต่ละจุด แต่บางปัญหาไม่ได้ต้องการรู้สนามทีละจุด เราต้องการรู้ว่าโดยรวมแล้วสนามไฟฟ้า ``ทะลุผ่าน'' พื้นผิวหนึ่งมากเพียงใด ปริมาณนี้เรียกว่า \textbf{ฟลักซ์ไฟฟ้า}

ให้พื้นผิวเล็ก ๆ มีพื้นที่ $dA$ และกำหนดเวกเตอร์พื้นที่ $d\vec{A}$ ให้มีขนาดเท่ากับ $dA$ และชี้ตั้งฉากออกจากพื้นผิว

สนามไฟฟ้าที่ผ่านพื้นผิวเล็กนี้ให้ฟลักซ์

\[
d\Phi_E=\vec{E}\cdot d\vec{A}
\]

ถ้าสนามสม่ำเสมอและพื้นผิวแบน มีพื้นที่ $A$ และมุมระหว่าง $\vec{E}$ กับเวกเตอร์ตั้งฉากของพื้นผิวเป็น $\theta$

\[
\Phi_E=EA\cos\theta
\]

สมการนี้มาจากการเลือกเฉพาะองค์ประกอบของสนามที่ตั้งฉากกับพื้นผิว

\[
E_\perp=E\cos\theta
\]

ดังนั้น

\[
\Phi_E=E_\perp A
\]

\begin{tcolorbox}[title=ความหมายเชิงภาพของฟลักซ์]
ฟลักซ์มากเมื่อสนามแรง พื้นที่ใหญ่ และสนามแทงผ่านพื้นผิวเกือบตั้งฉาก

ถ้าสนามขนานกับพื้นผิว จะมี $\theta=90^\circ$ และ

\[
\Phi_E=0
\]

เพราะไม่มีองค์ประกอบของสนามทะลุผ่านพื้นผิว
\end{tcolorbox}

\subsection{พื้นผิวปิดและทิศของเวกเตอร์พื้นที่}

สำหรับพื้นผิวปิด เวกเตอร์พื้นที่ทุกจุดนิยามให้ชี้ออกจากปริมาตรที่พื้นผิวล้อมไว้เสมอ ฟลักซ์รวมจึงเขียนเป็น

\[
\Phi_E=\oint_S \vec{E}\cdot d\vec{A}
\]

ฟลักซ์เป็นบวกเมื่อสนามชี้ออกจากพื้นผิว และเป็นลบเมื่อสนามชี้เข้า

\subsection{กฎของเกาส์}

กฎของเกาส์เป็นกฎพื้นฐานของไฟฟ้าสถิต ซึ่งเชื่อมฟลักซ์รวมผ่านพื้นผิวปิดกับประจุสุทธิที่อยู่ภายในพื้นผิวนั้น

\[
\oint_S \vec{E}\cdot d\vec{A}
=
\frac{Q_{\text{enc}}}{\varepsilon_0}
\]

\begin{tcolorbox}[title=กฎของเกาส์]
\[
\boxed{
\oint_S \vec{E}\cdot d\vec{A}
=
\frac{Q_{\text{enc}}}{\varepsilon_0}}
\]

$Q_{\text{enc}}$ คือประจุสุทธิที่อยู่ภายในพื้นผิวปิดเท่านั้น
\end{tcolorbox}

\subsection{พื้นผิวเกาส์คืออะไร}

\textbf{พื้นผิวเกาส์} คือพื้นผิวปิดสมมติที่เราเลือกขึ้นเพื่อใช้กฎของเกาส์ ไม่ใช่วัตถุจริงและไม่จำเป็นต้องมีประจุอยู่บนพื้นผิวนั้น

การเลือกพื้นผิวเกาส์ที่ดีต้องอาศัยสมมาตร เพื่อให้บนแต่ละส่วนของพื้นผิวเกิดกรณีใดกรณีหนึ่ง

\begin{enumerate}
    \item สนามมีขนาดคงที่และขนานกับ $d\vec{A}$ ทำให้ฟลักซ์เป็น $EA$
    \item สนามตั้งฉากกับ $d\vec{A}$ ทำให้ฟลักซ์เป็นศูนย์
    \item สนามเป็นศูนย์บนส่วนนั้น
\end{enumerate}

\begin{tcolorbox}[title=ข้อควรระวัง]
ประจุภายนอกพื้นผิวเกาส์ยังสร้างสนามไฟฟ้าบนพื้นผิวได้ แต่ฟลักซ์สุทธิที่เกิดจากประจุภายนอกผ่านพื้นผิวปิดรวมกันเป็นศูนย์ กฎของเกาส์จึงขึ้นกับประจุสุทธิภายในเท่านั้น
\end{tcolorbox}

\subsection{พิสูจน์สนามของแผ่นประจุขนาดใหญ่มาก}

พิจารณาแผ่นบางขนาดใหญ่มาก มีความหนาแน่นประจุผิวสม่ำเสมอ $\sigma$ เนื่องจากสมมาตร สนามไฟฟ้าต้องตั้งฉากกับแผ่น และมีขนาดเท่ากันทั้งสองด้าน

เลือกพื้นผิวเกาส์เป็นทรงกระบอกสั้นหรือ ``กล่องยา'' คร่อมแผ่น โดยหน้าตัดแต่ละด้านมีพื้นที่ $A$

ฟลักซ์ผ่านผิวด้านข้างเป็นศูนย์ เพราะสนามขนานกับผิวด้านข้าง จึงมี

\[
\vec{E}\cdot d\vec{A}=0
\]

ฟลักซ์ผ่านหน้าบนคือ

\[
\Phi_{\text{top}}=EA
\]

และฟลักซ์ผ่านหน้าล่างคือ

\[
\Phi_{\text{bottom}}=EA
\]

ดังนั้นฟลักซ์รวมเป็น

\[
\Phi_E=2EA
\]

ประจุที่พื้นผิวเกาส์ครอบไว้คือ

\[
Q_{\text{enc}}=\sigma A
\]

ใช้กฎของเกาส์

\[
2EA=\frac{\sigma A}{\varepsilon_0}
\]

ตัด $A$

\[
\boxed{E=\frac{\sigma}{2\varepsilon_0}}
\]

ถ้า $\sigma>0$ สนามชี้ออกจากแผ่นทั้งสองด้าน ถ้า $\sigma<0$ สนามชี้เข้าหาแผ่น

\subsection{สนามระหว่างแผ่นประจุสองแผ่น}

พิจารณาแผ่นขนานขนาดใหญ่มากสองแผ่น แผ่นหนึ่งมี $+\sigma$ และอีกแผ่นมี $-\sigma$

สนามจากแต่ละแผ่นมีขนาด

\[
\frac{\sigma}{2\varepsilon_0}
\]

บริเวณระหว่างแผ่น สนามจากทั้งสองแผ่นชี้จากแผ่นบวกไปหาแผ่นลบในทิศเดียวกัน จึงรวมกันเป็น

\[
E_{\text{between}}
=
\frac{\sigma}{2\varepsilon_0}
+
\frac{\sigma}{2\varepsilon_0}
\]

\[
\boxed{E_{\text{between}}=\frac{\sigma}{\varepsilon_0}}
\]

บริเวณภายนอกแผ่น สนามจากสองแผ่นมีขนาดเท่ากันแต่ทิศตรงข้าม จึงหักล้างกัน

\[
\boxed{E_{\text{outside}}\approx0}
\]

ผลนี้เป็นค่าประมาณที่ดีเมื่อขนาดของแผ่นใหญ่กว่าระยะห่างมาก และพิจารณาบริเวณที่ไม่ใกล้ขอบแผ่น

\begin{exercisebox}{แบบฝึกหัด: ฟลักซ์ไฟฟ้าและกฎของเกาส์}
\begin{enumerate}
    \item สนามสม่ำเสมอ $E$ ทำมุม $60^\circ$ กับเวกเตอร์ตั้งฉากของพื้นผิวพื้นที่ $A$ จงหาฟลักซ์ผ่านพื้นผิว
    \item เพราะเหตุใดฟลักซ์ผ่านผิวด้านข้างของพื้นผิวเกาส์รูปกล่องยาคร่อมแผ่นประจุจึงเป็นศูนย์
    \item แผ่นประจุขนาดใหญ่มากมีความหนาแน่นประจุผิว $-\sigma$ จงหาขนาดและทิศของสนามทั้งสองด้าน
    \item แผ่นขนานสองแผ่นมีความหนาแน่นประจุ $+\sigma$ และ $-\sigma$ จงอธิบายด้วยการซ้อนทับว่าทำไมสนามภายนอกจึงเป็นศูนย์โดยประมาณ
\end{enumerate}
\end{exercisebox}

% =========================================================
\section{งาน พลังงานศักย์ไฟฟ้า และศักย์ไฟฟ้า}

\subsection{งานของแรงไฟฟ้าและพลังงานศักย์ไฟฟ้า}

แรงไฟฟ้าสถิตเป็นแรงอนุรักษ์ ดังนั้นสามารถนิยามพลังงานศักย์ไฟฟ้าได้

สำหรับประจุจุดสองประจุที่อยู่ห่างกัน $r$

\[
U=k_e\frac{q_1q_2}{r}
\]

เครื่องหมายของ $U$ สำคัญมาก

ถ้า $q_1q_2>0$ จะได้ $U>0$

ถ้า $q_1q_2<0$ จะได้ $U<0$

\begin{tcolorbox}[title=พลังงานศักย์ไฟฟ้าของประจุสองตัว]
\[
U=k_e\frac{q_1q_2}{r}
\]

โดยเลือกให้ $U=0$ เมื่อประจุทั้งสองอยู่ห่างกันอนันต์
\end{tcolorbox}

ถ้ามีหลายประจุ พลังงานศักย์รวมของระบบคือผลรวมของทุกคู่

\[
U_{\text{total}}=\sum_{i<j} k_e\frac{q_iq_j}{r_{ij}}
\]

\subsection{ศักย์ไฟฟ้า}

ศักย์ไฟฟ้าคือพลังงานศักย์ไฟฟ้าต่อหนึ่งหน่วยประจุทดสอบ

\[
V=\frac{U}{q_0}
\]

ศักย์ไฟฟ้าจากประจุจุด $q$ ที่ระยะ $r$ คือ

\[
V=k_e\frac{q}{r}
\]

\begin{tcolorbox}[title=ศักย์ไฟฟ้าจากประจุจุด]
\[
V=k_e\frac{q}{r}
\]

ศักย์ไฟฟ้าเป็นสเกลาร์ จึงรวมกันด้วยเครื่องหมายบวกหรือลบตามค่าประจุ ไม่ต้องแตกเวกเตอร์
\end{tcolorbox}

\subsection{งานกับความต่างศักย์}

จากนิยาม

\[
U=qV
\]

จึงได้การเปลี่ยนพลังงานศักย์ไฟฟ้า

\[
\Delta U=q\Delta V
\]

งานที่แรงไฟฟ้าทำคือ

\[
W_{\text{electric}}=-\Delta U
\]

ดังนั้น

\[
W_{\text{electric}}=-q\Delta V
\]

\begin{tcolorbox}[title=ความสัมพันธ์งานและศักย์ไฟฟ้า]
\[
\Delta U=q\Delta V
\]

\[
W_{\text{electric}}=-\Delta U=-q\Delta V
\]
\end{tcolorbox}

\subsection{ความสัมพันธ์ระหว่างสนามไฟฟ้าและศักย์ไฟฟ้า}

จากนิยามของความต่างศักย์

\[
\Delta V=\frac{\Delta U}{q}
\]

และงานที่แรงไฟฟ้าทำสัมพันธ์กับพลังงานศักย์ด้วย

\[
W_{\text{electric}}=-\Delta U
\]

ถ้าประจุ $q$ เคลื่อนที่เป็นระยะเล็ก ๆ $d\vec{\ell}$ ในสนามไฟฟ้า งานเล็ก ๆ ที่แรงไฟฟ้าทำคือ

\[
dW_{\text{electric}}=\vec{F}\cdot d\vec{\ell}
\]

เพราะ

\[
\vec{F}=q\vec{E}
\]

จึงได้

\[
dW_{\text{electric}}=q\vec{E}\cdot d\vec{\ell}
\]

แต่

\[
dW_{\text{electric}}=-dU
\]

ดังนั้น

\[
dU=-q\vec{E}\cdot d\vec{\ell}
\]

หารด้วย $q$

\[
dV=-\vec{E}\cdot d\vec{\ell}
\]

อินทิเกรตจากจุด $i$ ไปจุด $f$

\[
\boxed{
V_f-V_i
=
-\int_i^f \vec{E}\cdot d\vec{\ell}}
\]

ถ้าสนามสม่ำเสมอและการกระจัดยาว $d$ อยู่ในทิศเดียวกับสนาม

\[
V_f-V_i=-Ed
\]

ดังนั้นขนาดของความต่างศักย์ระหว่างสองจุดคือ

\[
\boxed{|\Delta V|=Ed}
\]

หรือ

\[
\boxed{E=\frac{|\Delta V|}{d}}
\]

เครื่องหมายลบหมายความว่า เมื่อเคลื่อนที่ไปตามทิศของสนามไฟฟ้า ศักย์ไฟฟ้าจะลดลง

\begin{examplebox}{ตัวอย่าง: ประจุวิ่งเข้าหาประจุชนิดเดียวกัน}
ประจุ $+q$ มวล $m$ ถูกยิงจากระยะไกลมากเข้าหาประจุคงที่ $+Q$ ถ้าต้องการให้เข้าใกล้ที่สุดได้ระยะ $D$ ต้องมีพลังงานจลน์เริ่มต้นเท่าไร สมมติว่า $+Q$ ถูกตรึงอยู่กับที่

ที่ระยะไกลมาก เลือกให้พลังงานศักย์เป็นศูนย์

\[
U_i=0
\]

ที่จุดเข้าใกล้ที่สุด ความเร็วของ $+q$ เป็นศูนย์ชั่วขณะ

\[
K_f=0
\]

พลังงานศักย์ที่ระยะ $D$ คือ

\[
U_f=k_e\frac{qQ}{D}
\]

ใช้การอนุรักษ์พลังงาน

\[
K_i+U_i=K_f+U_f
\]

\[
K_i= k_e\frac{qQ}{D}
\]

ดังนั้น

\[
\boxed{K_i=k_e\frac{qQ}{D}}
\]
\end{examplebox}

\begin{exercisebox}{แบบฝึกหัด 3: ศักย์ไฟฟ้าและพลังงาน}
\begin{enumerate}
    \item ประจุ $+q$ และ $+2q$ อยู่ห่างกัน $a$ จงหาพลังงานศักย์ไฟฟ้าของระบบ
    \item ประจุ $+q$ และ $-q$ อยู่ห่างกัน $a$ จงหาพลังงานศักย์ไฟฟ้าของระบบ
    \item ประจุ $+Q$ อยู่ที่จุดกำเนิด จงหาศักย์ไฟฟ้าที่ระยะ $r$
    \item ประจุ $+q$ ถูกปล่อยจากหยุดนิ่งในบริเวณที่ศักย์ไฟฟ้าลดลงจาก $V_1$ เป็น $V_2$ จงหาการเปลี่ยนพลังงานจลน์ของประจุ
\end{enumerate}
\end{exercisebox}

% =========================================================
\section{ตัวนำและสมดุลไฟฟ้าสถิต}

\subsection{เหตุใดสนามภายในตัวนำจึงเป็นศูนย์}

ในตัวนำมีประจุอิสระที่เคลื่อนที่ได้ ถ้าภายในตัวนำยังมีสนามไฟฟ้า

\[
\vec{E}\neq0
\]

ประจุอิสระจะถูกแรง

\[
\vec{F}=q\vec{E}
\]

และยังคงเคลื่อนที่ต่อไป จึงยังไม่ใช่สมดุลไฟฟ้าสถิต ดังนั้นเมื่อระบบเข้าสู่สมดุลแล้ว ต้องมี

\[
\boxed{\vec{E}_{\text{inside}}=0}
\]

\subsection{เหตุใดประจุส่วนเกินจึงอยู่บนผิว}

สมมติว่ายังมีประจุสุทธิอยู่ภายในเนื้อตัวนำ เลือกพื้นผิวเกาส์ปิดที่อยู่ในเนื้อตัวนำทั้งหมด เพราะสนามภายในตัวนำเป็นศูนย์ ฟลักซ์ผ่านพื้นผิวนี้เป็นศูนย์

\[
\oint \vec{E}\cdot d\vec{A}=0
\]

ใช้กฎของเกาส์

\[
0=\frac{Q_{\text{enc}}}{\varepsilon_0}
\]

จึงได้

\[
Q_{\text{enc}}=0
\]

ดังนั้นในสมดุลไฟฟ้าสถิต ประจุส่วนเกินไม่อยู่ในปริมาตรของตัวนำ แต่กระจายอยู่บนผิว

\subsection{เหตุใดสนามที่ผิวต้องตั้งฉากกับผิว}

ถ้าสนามที่ผิวมีองค์ประกอบตามแนวสัมผัส

\[
E_{\parallel}\neq0
\]

องค์ประกอบนี้จะออกแรงให้ประจุอิสระเคลื่อนที่ไปตามผิว จึงยังไม่เกิดสมดุล ดังนั้น

\[
E_{\parallel}=0
\]

เหลือเฉพาะองค์ประกอบตั้งฉากกับผิว

\subsection{สนามไฟฟ้าทันทีภายนอกผิวตัวนำ}

ให้ผิวตัวนำมีความหนาแน่นประจุผิว $\sigma$ เลือกพื้นผิวเกาส์รูปกล่องยาสั้นมากคร่อมผิวตัวนำ โดยหน้าตัดมีพื้นที่ $A$

หน้าที่อยู่ภายในตัวนำไม่มีฟลักซ์ เพราะ

\[
E_{\text{inside}}=0
\]

ผิวด้านข้างไม่มีฟลักซ์ในขีดจำกัดที่กล่องบางมาก ส่วนหน้าด้านนอกมีฟลักซ์

\[
\Phi_E=EA
\]

ประจุที่ครอบไว้คือ

\[
Q_{\text{enc}}=\sigma A
\]

ใช้กฎของเกาส์

\[
EA=\frac{\sigma A}{\varepsilon_0}
\]

ตัด $A$

\[
\boxed{E_{\text{outside}}=\frac{\sigma}{\varepsilon_0}}
\]

ทิศตั้งฉากออกจากผิวเมื่อ $\sigma>0$ และตั้งฉากเข้าหาผิวเมื่อ $\sigma<0$

\subsection{ผิวตัวนำเป็นผิวศักย์เท่ากัน}

สำหรับจุดสองจุดบนหรือภายในตัวนำ

\[
V_f-V_i=-\int_i^f \vec{E}\cdot d\vec{\ell}
\]

แต่ในตัวนำและตามแนวสัมผัสผิวมีสนามเป็นศูนย์ จึงได้

\[
V_f-V_i=0
\]

ดังนั้นทุกจุดบนตัวนำที่สมดุลมีศักย์เท่ากัน

\subsection{ทรงกลมตัวนำ}

ให้ทรงกลมตัวนำรัศมี $R$ มีประจุรวม $Q$ สำหรับจุดภายนอกที่ระยะ $r\ge R$ เลือกพื้นผิวเกาส์ทรงกลมร่วมศูนย์กลางรัศมี $r$

จากสมมาตร สนามมีขนาดคงที่บนพื้นผิวและชี้ตามแนวรัศมี

\[
\oint \vec{E}\cdot d\vec{A}=E(4\pi r^2)
\]

ใช้กฎของเกาส์

\[
E(4\pi r^2)=\frac{Q}{\varepsilon_0}
\]

ดังนั้น

\[
E(r)=\frac{1}{4\pi\varepsilon_0}\frac{Q}{r^2}
=k_e\frac{Q}{r^2}
\]

หาศักย์ที่ผิวโดยเลือก $V(\infty)=0$

\[
V(R)-V(\infty)
=
-\int_{\infty}^{R}\vec{E}\cdot d\vec{\ell}
\]

\[
V(R)
=-\int_{\infty}^{R}k_e\frac{Q}{r^2}\,dr
\]

\[
V(R)=k_e\frac{Q}{R}
\]

เพราะสนามภายในตัวนำเป็นศูนย์ ศักย์ทุกจุดภายในจึงเท่ากับศักย์ที่ผิว

\[
\boxed{V_{\text{inside}}=V_{\text{surface}}=k_e\frac{Q}{R}}
\]

ถ้าทรงกลมตัวนำสองลูกถูกเชื่อมด้วยลวด ประจุจะไหลจนศักย์เท่ากัน

\[
k_e\frac{Q_1}{R_1}=k_e\frac{Q_2}{R_2}
\]

จึงได้

\[
\frac{Q_1}{R_1}=\frac{Q_2}{R_2}
\]

ร่วมกับการอนุรักษ์ประจุ

\[
Q_{\text{total}}=Q_1+Q_2
\]

% =========================================================
\section{ตัวเก็บประจุ}

\subsection{ความจุไฟฟ้า}

ตัวเก็บประจุประกอบด้วยตัวนำสองส่วนที่มีประจุขนาดเท่ากันแต่เครื่องหมายตรงข้าม

\[
+Q\qquad\text{และ}\qquad-Q
\]

เมื่อเพิ่มประจุ $Q$ ความต่างศักย์ระหว่างตัวนำเพิ่มตามสัดส่วน สำหรับรูปทรงและตัวกลางที่กำหนด อัตราส่วน $Q/\Delta V$ จึงคงที่

\begin{tcolorbox}[title=นิยามความจุไฟฟ้า]
\[
\boxed{C=\frac{Q}{\Delta V}}
\]

$Q$ หมายถึงขนาดประจุบนแผ่นใดแผ่นหนึ่ง และ $\Delta V$ หมายถึงขนาดความต่างศักย์ระหว่างตัวนำ
\end{tcolorbox}

ดังนั้น

\[
Q=C\Delta V
\]

หน่วยของความจุไฟฟ้าคือฟารัด

\[
1\,\si{F}=1\,\si{C/V}
\]

\subsection{พิสูจน์ความจุของตัวเก็บประจุแผ่นขนาน}

ให้แผ่นตัวนำขนานสองแผ่นมีพื้นที่ $A$ อยู่ห่างกัน $d$ และมีประจุ $+Q$ กับ $-Q$ สมมติว่าแผ่นมีขนาดใหญ่มากเมื่อเทียบกับ $d$ จึงละเลยสนามบริเวณขอบ

ความหนาแน่นประจุบนผิวด้านในคือ

\[
\sigma=\frac{Q}{A}
\]

จากกฎของเกาส์และการซ้อนทับของสนามจากแผ่น $+\sigma$ และ $-\sigma$ สนามระหว่างแผ่นมีขนาด

\[
E=\frac{\sigma}{\varepsilon_0}
\]

แทน $\sigma=Q/A$

\[
E=\frac{Q}{\varepsilon_0A}
\]

สนามระหว่างแผ่นสม่ำเสมอ ดังนั้นขนาดความต่างศักย์คือ

\[
\Delta V=Ed
\]

แทนค่า $E$

\[
\Delta V=\frac{Qd}{\varepsilon_0A}
\]

ใช้คำนิยาม

\[
C=\frac{Q}{\Delta V}
\]

จึงได้

\[
C
=
\frac{Q}{Qd/(\varepsilon_0A)}
\]

ตัด $Q$

\[
\boxed{C=\varepsilon_0\frac{A}{d}}
\]

\begin{tcolorbox}[title=ความหมายของผลลัพธ์]
\begin{enumerate}
    \item พื้นที่ $A$ มากขึ้น ทำให้เก็บประจุได้มากขึ้นที่ความต่างศักย์เดิม จึงทำให้ $C$ มากขึ้น
    \item ระยะห่าง $d$ มากขึ้น ทำให้ต้องใช้ประจุมากขึ้นเพื่อสร้างความต่างศักย์เท่าเดิม จึงทำให้ $C$ ลดลง
\end{enumerate}
\end{tcolorbox}

ถ้าระหว่างแผ่นมีไดอิเล็กทริกที่มีสภาพยอมไฟฟ้า $\varepsilon$ แทนสุญญากาศ จะได้

\[
\boxed{C=\varepsilon\frac{A}{d}}
\]

\subsection{พิสูจน์การต่อตัวเก็บประจุแบบขนาน}

เมื่อต่อขนาน ปลายทั้งสองของตัวเก็บประจุทุกตัวต่อกับจุดศักย์เดียวกัน จึงมีความต่างศักย์เท่ากัน

\[
\Delta V_1=\Delta V_2=\cdots=\Delta V
\]

ประจุบนแต่ละตัวคือ

\[
Q_i=C_i\Delta V
\]

ประจุรวมที่แหล่งกำเนิดจ่ายคือ

\[
Q_{\text{total}}=Q_1+Q_2+\cdots
\]

แทน $Q_i=C_i\Delta V$

\[
Q_{\text{total}}
=(C_1+C_2+\cdots)\Delta V
\]

แต่โดยนิยามของความจุสมมูล

\[
Q_{\text{total}}=C_{\text{eq}}\Delta V
\]

เปรียบเทียบสัมประสิทธิ์ของ $\Delta V$

\[
\boxed{C_{\text{eq}}=C_1+C_2+C_3+\cdots}
\]

\subsection{พิสูจน์การต่อตัวเก็บประจุแบบอนุกรม}

เมื่อต่ออนุกรม จุดเชื่อมภายในเป็นตัวนำที่แยกจากแหล่งกำเนิดและเดิมเป็นกลาง ประจุที่ถูกเหนี่ยวนำบนสองด้านของจุดเชื่อมจึงมีขนาดเท่ากันแต่เครื่องหมายตรงข้าม ทำให้ตัวเก็บประจุทุกตัวมีประจุขนาดเท่ากัน

\[
Q_1=Q_2=Q_3=Q
\]

ความต่างศักย์รวมคือ

\[
\Delta V_{\text{total}}
=
\Delta V_1+\Delta V_2+\Delta V_3+\cdots
\]

สำหรับแต่ละตัว

\[
\Delta V_i=\frac{Q}{C_i}
\]

ดังนั้น

\[
\Delta V_{\text{total}}
=
Q\left(
\frac{1}{C_1}
+
\frac{1}{C_2}
+
\frac{1}{C_3}
+
\cdots
\right)
\]

แต่สำหรับตัวเก็บประจุสมมูล

\[
\Delta V_{\text{total}}=\frac{Q}{C_{\text{eq}}}
\]

เปรียบเทียบแล้วได้

\[
\boxed{
\frac{1}{C_{\text{eq}}}
=
\frac{1}{C_1}
+
\frac{1}{C_2}
+
\frac{1}{C_3}
+
\cdots}
\]

\subsection{พิสูจน์พลังงานที่เก็บในตัวเก็บประจุ}

ขณะชาร์จตัวเก็บประจุ ความต่างศักย์ไม่ได้มีค่า $V$ ตั้งแต่เริ่ม แต่เพิ่มจากศูนย์ตามประจุที่สะสม

เมื่อมีประจุสะสมชั่วขณะเป็น $q$ ความต่างศักย์คือ

\[
v(q)=\frac{q}{C}
\]

การย้ายประจุเล็ก ๆ $dq$ จากแผ่นลบไปยังแผ่นบวกต้องใช้ผลงาน

\[
dW=v(q)\,dq
\]

ดังนั้นพลังงานรวมที่ใช้ชาร์จจาก $0$ ถึง $Q$ คือ

\[
U=\int_0^Q dW
\]

\[
U=\int_0^Q \frac{q}{C}\,dq
\]

\[
U=\frac{1}{C}\left[\frac{q^2}{2}\right]_0^Q
\]

\[
U=\frac{Q^2}{2C}
\]

ใช้ $Q=CV$ จะได้รูปเทียบเท่า

\[
\boxed{
U=\frac{Q^2}{2C}
=\frac{1}{2}QV
=\frac{1}{2}CV^2}
\]

ตัวประกอบ $1/2$ เกิดขึ้นเพราะความต่างศักย์เฉลี่ยระหว่างการชาร์จมีค่า $V/2$ ไม่ใช่ $V$

\subsection{พิสูจน์พลังงานต่อหน่วยปริมาตรในสนามไฟฟ้า}

สำหรับตัวเก็บประจุแผ่นขนาน

\[
U=\frac{1}{2}CV^2
\]

แทน

\[
C=\varepsilon_0\frac{A}{d}
\]

และ

\[
V=Ed
\]

จะได้

\[
U
=
\frac{1}{2}
\left(\varepsilon_0\frac{A}{d}\right)
(Ed)^2
\]

\[
U=\frac{1}{2}\varepsilon_0E^2Ad
\]

ปริมาตรบริเวณที่มีสนามระหว่างแผ่นคือ

\[
\mathcal{V}=Ad
\]

ดังนั้นพลังงานต่อหน่วยปริมาตรคือ

\[
u_E=\frac{U}{\mathcal{V}}
\]

\[
\boxed{u_E=\frac{1}{2}\varepsilon_0E^2}
\]

ผลนี้แสดงว่าพลังงานไม่ได้อยู่เพียง ``บนแผ่น'' แต่สามารถมองว่าเก็บอยู่ในสนามไฟฟ้าระหว่างแผ่น

\subsection{พิสูจน์แรงดึงดูดระหว่างแผ่น}

จุดที่ต้องระวังคือ แผ่นหนึ่งไม่ออกแรงต่อตัวเอง ดังนั้นแรงบนแผ่นบวกต้องคำนวณจากสนามที่สร้างโดยแผ่นลบเพียงแผ่นเดียว

สนามจากแผ่นลบเพียงแผ่นเดียวมีขนาด

\[
E_{\text{other}}=\frac{\sigma}{2\varepsilon_0}
\]

ประจุบนแผ่นบวกคือ

\[
Q=\sigma A
\]

ดังนั้นแรงดึงดูดมีขนาด

\[
F=QE_{\text{other}}
\]

\[
F=(\sigma A)\left(\frac{\sigma}{2\varepsilon_0}\right)
\]

\[
\boxed{F=\frac{\sigma^2A}{2\varepsilon_0}}
\]

เพราะสนามรวมระหว่างแผ่นคือ

\[
E=\frac{\sigma}{\varepsilon_0}
\]

จึงเขียนได้เป็น

\[
\boxed{F=\frac{1}{2}\varepsilon_0E^2A}
\]

ถ้าแผ่นต่อกับแหล่งกำเนิดที่รักษาความต่างศักย์ $V$ คงที่

\[
E=\frac{V}{d}
\]

ดังนั้น

\[
F
=
\frac{1}{2}\varepsilon_0A\left(\frac{V}{d}\right)^2
\]

และเพราะ

\[
C=\varepsilon_0\frac{A}{d}
\]

จึงได้

\[
\boxed{F=\frac{CV^2}{2d}}
\]

\begin{tcolorbox}[title=เหตุผลของตัวประกอบ $1/2$]
ห้ามใช้สนามรวมระหว่างแผ่น $E=\sigma/\varepsilon_0$ คูณกับประจุบนแผ่นโดยตรง เพราะสนามรวมนั้นรวมสนามที่แผ่นสร้างเองด้วย สนามที่ทำให้แผ่นหนึ่งถูกดึงดูดคือสนามจากอีกแผ่นเท่านั้น จึงมีขนาดครึ่งหนึ่งของสนามรวม
\end{tcolorbox}

\begin{examplebox}{ตัวอย่าง: ตัวเก็บประจุสองตัวต่ออนุกรม}
ตัวเก็บประจุ $C_1=2\,\si{\micro F}$ และ $C_2=6\,\si{\micro F}$ ต่ออนุกรมกับแบตเตอรี่ $12\,\si{V}$ จงหาประจุบนตัวเก็บประจุแต่ละตัว

หาความจุรวม

\[
\frac{1}{C_{\text{eq}}}
=
\frac{1}{2}+\frac{1}{6}
=
\frac{4}{6}
\]

\[
C_{\text{eq}}=\frac{3}{2}\,\si{\micro F}
\]

ประจุที่แหล่งกำเนิดจ่ายคือ

\[
Q=C_{\text{eq}}V
\]

\[
Q=\left(\frac{3}{2}\right)(12)
=18\,\si{\micro C}
\]

เพราะต่ออนุกรม ประจุบนแต่ละตัวมีขนาดเท่ากัน

\[
\boxed{Q_1=Q_2=18\,\si{\micro C}}
\]
\end{examplebox}

\begin{exercisebox}{แบบฝึกหัด 4: ตัวเก็บประจุ}
\begin{enumerate}
    \item ใช้กฎของเกาส์พิสูจน์ว่าสนามระหว่างแผ่นขนานขนาดใหญ่มากที่มี $+\sigma$ และ $-\sigma$ มีขนาด $\sigma/\varepsilon_0$
    \item จากผลข้อก่อนหน้า พิสูจน์ว่า $C=\varepsilon_0A/d$
    \item ตัวเก็บประจุ $2\,\si{\micro F}$ และ $3\,\si{\micro F}$ ต่อขนานกัน จงหาความจุรวมโดยเริ่มจากการอนุรักษ์ประจุ
    \item ตัวเก็บประจุ $2\,\si{\micro F}$ และ $3\,\si{\micro F}$ ต่ออนุกรมกัน จงหาความจุรวมโดยเริ่มจากผลรวมความต่างศักย์
    \item พิสูจน์ว่า $U=\frac{1}{2}CV^2$ โดยอินทิเกรตงานระหว่างการชาร์จ
    \item ใช้สนามจากแผ่นตรงข้ามพิสูจน์ว่าแรงดึงดูดระหว่างแผ่นคือ $F=\frac{1}{2}\varepsilon_0E^2A$
\end{enumerate}
\end{exercisebox}

% =========================================================
\section{วิธีเลือกหลักในการแก้โจทย์}

\begin{tcolorbox}[title=แผนแก้โจทย์ไฟฟ้าสถิต]
\begin{enumerate}
    \item ถ้าถามแรง ให้เริ่มจากกฎของคูลอมบ์ และรวมแบบเวกเตอร์
    \item ถ้าถามสนามไฟฟ้า ให้หา $\vec{E}$ จากแต่ละประจุ แล้วรวมแบบเวกเตอร์
    \item ถ้าถามศักย์ไฟฟ้า ให้หา $V$ จากแต่ละประจุ แล้วรวมแบบสเกลาร์
    \item ถ้าถามพลังงานของระบบหลายประจุ ให้รวมพลังงานของทุกคู่
    \item ถ้ามีสมดุล ให้ใช้ $\sum \vec{F}=0$ ร่วมกับแรงไฟฟ้า
    \item ถ้ามีการเคลื่อนที่เข้าหาหรือหนีจากประจุ ให้ใช้พลังงานกลรวมร่วมกับ $U=kq_1q_2/r$
    \item ถ้าเป็นตัวเก็บประจุ ให้ดูว่าอนุกรมหรือขนานก่อน แล้วใช้ $Q=CV$ และ $U=\frac{1}{2}CV^2$
\end{enumerate}
\end{tcolorbox}

% =========================================================
\section{ชุดโจทย์รวมท้ายบท: ระดับ สอวน. ค่าย 1}

\subsection*{ระดับพื้นฐาน}

\begin{enumerate}
    \item ประจุ $+2.0\,\si{\micro C}$ และ $-3.0\,\si{\micro C}$ อยู่ห่างกัน $0.50\,\si{m}$ จงหาขนาดแรงไฟฟ้าระหว่างกัน
    \item ประจุ $+Q$ อยู่ที่จุดกำเนิด จงหาสนามไฟฟ้าและศักย์ไฟฟ้าที่ระยะ $r$
    \item ตัวเก็บประจุ $5\,\si{\micro F}$ ต่อกับแบตเตอรี่ $20\,\si{V}$ จงหา $Q$ และ $U$
    \item ประจุ $+q$ และ $-q$ อยู่ห่างกัน $a$ จงหาพลังงานศักย์ไฟฟ้าของระบบ
\end{enumerate}

\subsection*{ระดับกลาง}

\begin{enumerate}
    \setcounter{enumi}{4}
    \item ประจุ $+q$ สองตัวอยู่ที่ $(0,a)$ และ $(0,-a)$ จงหาขนาดสนามไฟฟ้าที่จุด $(x,0)$
    \item วงแหวนประจุรัศมี $R$ มีประจุรวม $Q$ จงหาตำแหน่งบนแกนที่สนามไฟฟ้ามีขนาดมากที่สุด
    \item ประจุ $+q$ มวล $m$ ถูกยิงจากระยะไกลมากเข้าหาประจุ $+Q$ ที่ถูกตรึงไว้ ถ้าพลังงานจลน์เริ่มต้นคือ $K$ จงหาระยะเข้าใกล้ที่สุด
    \item ทรงกลมตัวนำรัศมี $a$ และ $b$ มีประจุเริ่มต้นเท่ากันลูกละ $Q$ แล้วเชื่อมด้วยลวด จงหาประจุสุดท้ายบนแต่ละลูก
\end{enumerate}

\subsection*{ระดับยาก}

\begin{enumerate}
    \setcounter{enumi}{8}
    \item ประจุ $+q$ สามตัวอยู่ที่มุมของสามเหลี่ยมด้านเท่าด้านยาว $a$ จงหาพลังงานศักย์ไฟฟ้ารวมของระบบ และขนาดแรงลัพธ์บนประจุหนึ่งตัว
    \item ประจุ $+q$ มวล $m$ ถูกยิงจากระยะไกลมากเข้าหาประจุ $+Q$ มวล $M$ ที่เดิมอยู่นิ่งและไม่ถูกตรึง ถ้าระยะเข้าใกล้ที่สุดคือ $D$ จงหาพลังงานจลน์เริ่มต้นที่ต้องใช้
    \item ตัวเก็บประจุแผ่นขนานมีความจุ $C$ ต่อกับแหล่งกำเนิดแรงเคลื่อนไฟฟ้า $\mathcal{E}$ แผ่นห่างกัน $d$ จงหาแรงดึงดูดระหว่างแผ่นโดยใช้แนวคิดพลังงาน
    \item ลวดครึ่งวงกลมรัศมี $R$ มีประจุกระจายสม่ำเสมอด้วยความหนาแน่นเชิงเส้น $\lambda$ จงหาศักย์ไฟฟ้าที่ศูนย์กลาง
\end{enumerate}

% =========================================================
\section{เฉลยย่อท้ายบท}

\begin{enumerate}
    \item
    \[
    F=k_e\frac{|q_1q_2|}{r^2}
    =k_e\frac{(2.0\times10^{-6})(3.0\times10^{-6})}{(0.50)^2}
    \]
    เป็นแรงดึงดูด

    \item
    \[
    E=k_e\frac{Q}{r^2}
    \]
    ชี้ออกจาก $+Q$
    \[
    V=k_e\frac{Q}{r}
    \]

    \item
    \[
    Q=CV=(5)(20)=100\,\si{\micro C}
    \]
    \[
    U=\frac{1}{2}CV^2=\frac{1}{2}(5\times10^{-6})(20)^2=1.0\times10^{-3}\,\si{J}
    \]

    \item
    \[
    U=k_e\frac{(+q)(-q)}{a}=-k_e\frac{q^2}{a}
    \]

    \item
    ระยะจากแต่ละประจุถึงจุดคือ
    \[
    r=\sqrt{x^2+a^2}
    \]
    องค์ประกอบแนว $y$ หักล้างกัน เหลือแนว $x$
    \[
    E=2\left(k_e\frac{q}{r^2}\right)\frac{x}{r}
    \]
    \[
    E=\frac{2k_eqx}{(x^2+a^2)^{3/2}}
    \]

    \item
    \[
    E(x)=k_e\frac{Qx}{(R^2+x^2)^{3/2}}
    \]
    หา $dE/dx=0$ จะได้
    \[
    x=\frac{R}{\sqrt{2}}
    \]

    \item
    ที่ระยะเข้าใกล้ที่สุด ความเร็วเป็นศูนย์
    \[
    K=k_e\frac{qQ}{r_{\min}}
    \]
    ดังนั้น
    \[
    r_{\min}=k_e\frac{qQ}{K}
    \]

    \item
    ศักย์เท่ากันหลังเชื่อม
    \[
    k_e\frac{Q_a}{a}=k_e\frac{Q_b}{b}
    \]
    และ
    \[
    Q_a+Q_b=2Q
    \]
    ดังนั้น
    \[
    Q_a=\frac{2Qa}{a+b}
    \]
    \[
    Q_b=\frac{2Qb}{a+b}
    \]

    \item
    พลังงานศักย์รวมมี 3 คู่
    \[
    U=3k_e\frac{q^2}{a}
    \]
    แรงจากอีกสองประจุทำมุมกัน $60^\circ$
    \[
    F_{\text{net}}=\sqrt{F^2+F^2+2F^2\cos60^\circ}
    \]
    \[
    F_{\text{net}}=\sqrt{3}F=\sqrt{3}k_e\frac{q^2}{a^2}
    \]

    \item
    ให้พลังงานจลน์เริ่มต้นเป็น $K$

    โมเมนตัมรวมคงตัว ที่ระยะใกล้ที่สุดความเร็วสัมพัทธ์เป็นศูนย์ ทั้งสองประจุมีความเร็วร่วมกัน
    \[
    V=\frac{m}{m+M}u
    \]
    ถ้า
    \[
    K=\frac{1}{2}mu^2
    \]
    พลังงานจลน์ของการเคลื่อนที่ของจุดศูนย์กลางมวลที่จุดใกล้ที่สุดคือ
    \[
    K_{\text{cm}}=\frac{m}{m+M}K
    \]
    ใช้พลังงาน
    \[
    K=K_{\text{cm}}+k_e\frac{qQ}{D}
    \]
    \[
    K\left(1-\frac{m}{m+M}\right)=k_e\frac{qQ}{D}
    \]
    \[
    K\frac{M}{m+M}=k_e\frac{qQ}{D}
    \]
    ดังนั้น
    \[
    K=\frac{m+M}{M}k_e\frac{qQ}{D}
    \]

    \item
    สำหรับแผ่นขนานที่ต่อกับแหล่งกำเนิด $\mathcal{E}$
    \[
    U=\frac{1}{2}C\mathcal{E}^2
    \]
    และ $C\propto 1/d$ จึงได้แรงดึงดูด
    \[
    F=\frac{C\mathcal{E}^2}{2d}
    \]

    \item
    ทุกชิ้นประจุอยู่ห่างจากศูนย์กลางเท่ากับ $R$
    \[
    V=k_e\int\frac{dq}{R}
    \]
    \[
    V=k_e\frac{Q}{R}
    \]
    สำหรับครึ่งวงกลม
    \[
    Q=\lambda\pi R
    \]
    ดังนั้น
    \[
    V=k_e\lambda\pi=\frac{\lambda}{4\varepsilon_0}
    \]
\end{enumerate}

% =========================================================
\section{สรุปสูตรสำคัญ}

\begin{tcolorbox}[title=สรุปบท: ไฟฟ้าสถิต]
กฎของคูลอมบ์

\[
F=k_e\frac{|q_1q_2|}{r^2}
\]

\[
k_e=\frac{1}{4\pi\varepsilon_0}
\]

สนามไฟฟ้า

\[
\vec{E}=\frac{\vec{F}}{q_0}
\]

\[
E=k_e\frac{|q|}{r^2}
\]

\[
\vec{E}_{\text{net}}=\vec{E}_1+\vec{E}_2+\cdots
\]

ฟลักซ์ไฟฟ้าและกฎของเกาส์

\[
\Phi_E=\oint_S\vec{E}\cdot d\vec{A}
\]

\[
\oint_S\vec{E}\cdot d\vec{A}=\frac{Q_{\text{enc}}}{\varepsilon_0}
\]

แผ่นประจุขนาดใหญ่มาก

\[
E=\frac{\sigma}{2\varepsilon_0}
\]

ระหว่างแผ่น $+\sigma$ และ $-\sigma$

\[
E=\frac{\sigma}{\varepsilon_0}
\]

สนามบนแกนของวงแหวนประจุ

\[
E_x=k_e\frac{Qx}{(R^2+x^2)^{3/2}}
\]

พลังงานศักย์ไฟฟ้า

\[
U=k_e\frac{q_1q_2}{r}
\]

\[
U_{\text{total}}=\sum_{i<j}k_e\frac{q_iq_j}{r_{ij}}
\]
\end{tcolorbox}

\begin{tcolorbox}[title=สรุปบท: ไฟฟ้าสถิต ต่อ]
ศักย์ไฟฟ้า

\[
V=\frac{U}{q_0}
\]

\[
V=k_e\frac{q}{r}
\]

\[
\Delta U=q\Delta V
\]

\[
W_{\text{electric}}=-\Delta U
\]

ตัวเก็บประจุ

\[
C=\frac{Q}{V}
\]

\[
C=\varepsilon_0\frac{A}{d}
\]

ต่อขนาน

\[
C_{\text{eq}}=C_1+C_2+\cdots
\]

ต่ออนุกรม

\[
\frac{1}{C_{\text{eq}}}=\frac{1}{C_1}+\frac{1}{C_2}+\cdots
\]

พลังงานในตัวเก็บประจุ

\[
U=\frac{1}{2}QV=\frac{1}{2}CV^2=\frac{Q^2}{2C}
\]

พลังงานต่อหน่วยปริมาตรในสนามไฟฟ้า

\[
u_E=\frac{1}{2}\varepsilon_0E^2
\]

แรงระหว่างแผ่นตัวเก็บประจุที่ต่อกับแหล่งจ่าย

\[
F=\frac{CV^2}{2d}
\]
\end{tcolorbox}

% =========================================================
\section{ข้อสอบจริง สอวน. ที่ใช้ความรู้บทไฟฟ้าสถิต}



% ---------------------------------------------------------
\subsection{ปี 2560 ข้อ 12: สนามไฟฟ้าจากประจุสองตัวที่จุด $(3,4)$}

\vspace{1em}

\IfFileExists{img/posn60q12.png}{
\begin{center}
    \includegraphics[width=1\linewidth]{img/posn60q12.png}
\end{center}
}{
\begin{tcolorbox}[title=ตำแหน่งภาพที่แนะนำ]
ให้ตัดภาพข้อสอบปี 2560 ข้อ 12 แล้วบันทึกเป็น
\[
\texttt{img/posn60q12.png}
\]
\end{tcolorbox}
}

\begin{examplebox}{เฉลยแนวคิด}
จุดที่โจทย์ให้คือ $(3,4)$ ซึ่งอยู่ห่างจากประจุที่ $(0,0)$ และประจุที่ $(6,0)$ เท่ากัน

\[
r=\sqrt{3^2+4^2}=5
\]

ดังนั้นสนามจากประจุทั้งสองมีขนาดเท่ากัน ถ้าประจุทั้งสองมีขนาดเท่ากันและเครื่องหมายตรงข้าม องค์ประกอบแนวดิ่งของสนามจะหักล้างกัน และองค์ประกอบแนวราบจะรวมกัน

แนวคิดที่ต้องใช้คือ

\[
\vec{E}_{\text{net}}=\vec{E}_1+\vec{E}_2
\]

และแตกเวกเตอร์สนามแต่ละตัวเป็นองค์ประกอบตามแกน $x$ และ $y$
\end{examplebox}

\begin{examplebox}{เฉลยละเอียด}
ทรงกลมตัวนำสองลูกมีรัศมี $a$ และ $b$ โดยเริ่มต้นมีประจุลูกละ $+Q$

ดังนั้นประจุรวมก่อนเชื่อมคือ

\[
Q_{\text{total}}=Q+Q=2Q
\]

เมื่อเชื่อมทรงกลมทั้งสองด้วยลวด ตัวนำทั้งระบบเข้าสู่สมดุลไฟฟ้าสถิต ประจุสามารถถ่ายเทระหว่างทรงกลมได้ แต่ประจุรวมยังคงเดิม

โจทย์กำหนดว่าหลังเชื่อม ทรงกลมรัศมี $a$ มีประจุ

\[
q_a=\frac{Q}{3}
\]

ดังนั้นจากการอนุรักษ์ประจุ

\[
q_a+q_b=2Q
\]

แทนค่า $q_a=Q/3$

\[
\frac{Q}{3}+q_b=2Q
\]

\[
q_b=2Q-\frac{Q}{3}
\]

\[
q_b=\frac{5Q}{3}
\]

เมื่อทรงกลมตัวนำทั้งสองเชื่อมด้วยลวดและอยู่ในสมดุล ศักย์ไฟฟ้าของทรงกลมทั้งสองต้องเท่ากัน

\[
V_a=V_b
\]

สำหรับทรงกลมตัวนำที่อยู่ห่างกันมากพอจนละเลยอิทธิพลระหว่างกันได้ ศักย์ที่ผิวทรงกลมคือ

\[
V=k_e\frac{q}{R}
\]

ดังนั้น

\[
k_e\frac{q_a}{a}
=
k_e\frac{q_b}{b}
\]

ตัด $k_e$

\[
\frac{q_a}{a}
=
\frac{q_b}{b}
\]

แทนค่า

\[
\frac{Q/3}{a}
=
\frac{5Q/3}{b}
\]

ตัด $Q/3$

\[
\frac{1}{a}
=
\frac{5}{b}
\]

จึงได้

\[
b=5a
\]

ดังนั้น

\[
\boxed{\frac{b}{a}=5}
\]
\end{examplebox}
% ---------------------------------------------------------
\subsection{ปี 2561 ข้อ 12: ประจุแขวนด้วยเชือกและสมดุลแรง}

\vspace{1em}

\IfFileExists{img/posn61q12.png}{
\begin{center}
    \includegraphics[width=1\linewidth]{img/posn61q12.png}
\end{center}
}{
\begin{tcolorbox}[title=ตำแหน่งภาพที่แนะนำ]
ให้ตัดภาพข้อสอบปี 2561 ข้อ 12 แล้วบันทึกเป็น
\[
\texttt{img/posn61q12.png}
\]
\end{tcolorbox}
}

\begin{examplebox}{เฉลยละเอียด}
ในสมดุลเดิม มุมกางของเชือกคงที่ จึงได้ว่าอัตราส่วนระหว่างแรงไฟฟ้ากับน้ำหนักคงที่

\[
\tan\theta=\frac{F_e}{mg}
\]

เพราะมุมเท่าเดิม ระยะห่างระหว่างประจุเท่าเดิมด้วย ดังนั้น

\[
\frac{F_e}{mg}=\text{constant}
\]

แรงไฟฟ้ามีขนาด

\[
F_e=k_e\frac{q^2}{r^2}
\]

เมื่อมวลเปลี่ยนจาก $m$ เป็น $2m$ และต้องการให้มุมเท่าเดิม ต้องมี

\[
\frac{k_e q'^2/r^2}{2mg}=\frac{k_e q^2/r^2}{mg}
\]

ตัดค่าคงที่ร่วมกัน

\[
\frac{q'^2}{2}=q^2
\]

\[
q'^2=2q^2
\]

ดังนั้น

\[
\boxed{q'=\sqrt{2}\,q}
\]
\end{examplebox}

% ---------------------------------------------------------
\subsection{ปี 2561 ข้อ 14: ตัวเก็บประจุในวงจรผสม}

\vspace{1em}

\IfFileExists{img/posn61q14.png}{
\begin{center}
    \includegraphics[width=1\linewidth]{img/posn61q14.png}
\end{center}
}{
\begin{tcolorbox}[title=ตำแหน่งภาพที่แนะนำ]
ให้ตัดภาพข้อสอบปี 2561 ข้อ 14 แล้วบันทึกเป็น
\[
\texttt{img/posn61q14.png}
\]
\end{tcolorbox}
}

\begin{examplebox}{เฉลยแนวคิด}
ข้อนี้ใช้การลดวงจรตัวเก็บประจุทีละชั้น

\begin{enumerate}
    \item รวมตัวที่ต่อขนานด้วย
    \[
    C_{\text{eq}}=C_1+C_2+\cdots
    \]
    \item รวมตัวที่ต่ออนุกรมด้วย
    \[
    \frac{1}{C_{\text{eq}}}=\frac{1}{C_1}+\frac{1}{C_2}+\cdots
    \]
    \item เมื่อได้ความจุรวม ใช้
    \[
    Q_{\text{total}}=C_{\text{eq}}V
    \]
    \item ย้อนกลับไปหาประจุของตัวเก็บประจุที่โจทย์ถาม โดยจำว่าอนุกรมมีประจุเท่ากัน ส่วนขนานมีความต่างศักย์เท่ากัน
\end{enumerate}
\end{examplebox}

% ---------------------------------------------------------
\subsection{ปี 2562 ข้อ 9: ศักย์ไฟฟ้าที่ศูนย์กลางของลวดครึ่งวงกลม}

\vspace{1em}

\IfFileExists{img/posn62q9.png}{
\begin{center}
    \includegraphics[width=1\linewidth]{img/posn62q9.png}
\end{center}
}{
\begin{tcolorbox}[title=ตำแหน่งภาพที่แนะนำ]
ให้ตัดภาพข้อสอบปี 2562 ข้อ 9 แล้วบันทึกเป็น
\[
\texttt{img/posn62q9.png}
\]
\end{tcolorbox}
}

\begin{examplebox}{เฉลยละเอียด}
ทุกชิ้นประจุบนลวดครึ่งวงกลมอยู่ห่างจากจุดศูนย์กลางเท่ากันคือ $R$

ศักย์ไฟฟ้าเป็นสเกลาร์ ดังนั้นรวมได้โดยตรง

\[
V=k_e\int\frac{dq}{R}
\]

เพราะ $R$ คงที่

\[
V=k_e\frac{1}{R}\int dq
\]

ประจุรวมบนครึ่งวงกลมคือ

\[
Q=\lambda(\pi R)
\]

ดังนั้น

\[
V=k_e\frac{\lambda\pi R}{R}
\]

\[
V=k_e\lambda\pi
\]

แทน $k_e=1/(4\pi\varepsilon_0)$

\[
V=\frac{1}{4\pi\varepsilon_0}\lambda\pi
\]

\[
\boxed{V=\frac{\lambda}{4\varepsilon_0}}
\]
\end{examplebox}

% ---------------------------------------------------------
\subsection{ปี 2562 ข้อ 25: สนามไฟฟ้าบนแกนของวงแหวนประจุ}

\vspace{1em}

\IfFileExists{img/posn62q25.png}{
\begin{center}
    \includegraphics[width=1\linewidth]{img/posn62q25.png}
\end{center}
}{
\begin{tcolorbox}[title=ตำแหน่งภาพที่แนะนำ]
ให้ตัดภาพข้อสอบปี 2562 ข้อ 25 แล้วบันทึกเป็น
\[
\texttt{img/posn62q25.png}
\]
\end{tcolorbox}
}

\begin{examplebox}{เฉลยละเอียด}
วงแหวนรัศมี $R$ มีประจุกระจายสม่ำเสมอด้วยความหนาแน่นเชิงเส้น $\lambda$ ดังนั้นประจุรวมคือ

\[
Q=2\pi R\lambda
\]

สนามไฟฟ้าบนแกนของวงแหวนที่ห่างจากศูนย์กลางเป็นระยะ $x$ คือ

\[
E=k_e\frac{Qx}{(R^2+x^2)^{3/2}}
\]

แทน $Q=2\pi R\lambda$

\[
E=\frac{1}{4\pi\varepsilon_0}\frac{(2\pi R\lambda)x}{(R^2+x^2)^{3/2}}
\]

จึงได้

\[
\boxed{E=\frac{\lambda R x}{2\varepsilon_0(R^2+x^2)^{3/2}}}
\]

ทิศของสนามอยู่ตามแกนของวงแหวน และถ้าประจุเป็นบวกกับ $x>0$ สนามชี้ไปทางแกนบวก
\end{examplebox}

% ---------------------------------------------------------
\subsection{อ่านเสริม: ข้อสอบศูนย์ สอวน. ปี 2563 ข้อ 6 พลังงานศักย์ไฟฟ้าของระบบประจุ}

\vspace{1em}

\IfFileExists{img/posn63centerq6.png}{
\begin{center}
    \includegraphics[width=1\linewidth]{img/posn63centerq6.png}
\end{center}
}{
\begin{tcolorbox}[title=ตำแหน่งภาพที่แนะนำ]
ให้ตัดภาพข้อสอบศูนย์ สอวน. ปี 2563 ข้อ 6 แล้วบันทึกเป็น
\[
\texttt{img/posn63centerq6.png}
\]
\end{tcolorbox}
}

\begin{examplebox}{เฉลยแนวคิด}
ข้อนี้ใช้หลักเดียวคือพลังงานศักย์ไฟฟ้ารวมของระบบหลายประจุ ต้องรวมทุกคู่เพียงครั้งเดียว

\[
U_{\text{total}}=\sum_{i<j}k_e\frac{q_iq_j}{r_{ij}}
\]

ขั้นตอนคือ

\begin{enumerate}
    \item ระบุประจุทุกตัวในรูป
    \item เขียนคู่ประจุทั้งหมด
    \item หาระยะระหว่างคู่ประจุแต่ละคู่
    \item ใส่เครื่องหมายของประจุใน $q_iq_j$ โดยตรง
    \item รวมทุกพจน์แบบสเกลาร์
\end{enumerate}
\end{examplebox}

% ---------------------------------------------------------
\subsection{อ่านเสริม: ข้อสอบศูนย์ สอวน. ปี 2563 ข้อ 7 สนามไฟฟ้าบนแกนของวงแหวน}

\vspace{1em}

\IfFileExists{img/posn63centerq7.png}{
\begin{center}
    \includegraphics[width=1\linewidth]{img/posn63centerq7.png}
\end{center}
}{
\begin{tcolorbox}[title=ตำแหน่งภาพที่แนะนำ]
ให้ตัดภาพข้อสอบศูนย์ สอวน. ปี 2563 ข้อ 7 แล้วบันทึกเป็น
\[
\texttt{img/posn63centerq7.png}
\]
\end{tcolorbox}
}

\begin{examplebox}{เฉลยละเอียด}
วงแหวนรัศมี $R$ มีประจุรวม $+q$ กระจายสม่ำเสมอ จุด $P$ อยู่บนแกนของวงแหวนห่างจากศูนย์กลางเป็นระยะ $x$

จากสมมาตร องค์ประกอบของสนามที่ตั้งฉากกับแกนหักล้างกันหมด เหลือเฉพาะองค์ประกอบตามแกน

\[
E=k_e\frac{qx}{(R^2+x^2)^{3/2}}
\]

ดังนั้น

\[
\boxed{E=\frac{1}{4\pi\varepsilon_0}\frac{qx}{(R^2+x^2)^{3/2}}}
\]

ทิศชี้ออกจากวงแหวนตามแกนไปยังจุด $P$ ถ้า $q>0$
\end{examplebox}

% ---------------------------------------------------------
\subsection{ปี 2566 ข้อ 12: กราฟสนามไฟฟ้าบนแกนของวงแหวนประจุ}

\vspace{1em}

\IfFileExists{img/posn66q12.png}{
\begin{center}
    \includegraphics[width=1\linewidth]{img/posn66q12.png}
\end{center}
}{
\begin{tcolorbox}[title=ตำแหน่งภาพที่แนะนำ]
ให้ตัดภาพข้อสอบปี 2566 ข้อ 12 แล้วบันทึกเป็น
\[
\texttt{img/posn66q12.png}
\]
\end{tcolorbox}
}

\begin{examplebox}{เฉลยละเอียด}
สำหรับวงแหวนประจุบวกรัศมี $R$ สนามบนแกน $x$ คือ

\[
E(x)=k_e\frac{Qx}{(R^2+x^2)^{3/2}}
\]

ตรวจลักษณะกราฟ

ที่ $x=0$

\[
E(0)=0
\]

เมื่อ $x>0$ เล็ก ๆ สนามเป็นบวกและเพิ่มขึ้น

เมื่อ $x$ ใหญ่มาก

\[
E(x)\approx k_e\frac{Q}{x^2}
\]

จึงลดลงเข้าใกล้ศูนย์

หาจุดที่มากที่สุดจาก

\[
\frac{d}{dx}\left[\frac{x}{(R^2+x^2)^{3/2}}\right]=0
\]

จะได้

\[
x=\frac{R}{\sqrt{2}}
\]

ดังนั้นกราฟต้องเริ่มที่ศูนย์ เพิ่มขึ้นถึงค่าสูงสุด แล้วลดลงเข้าใกล้ศูนย์
\end{examplebox}

% ---------------------------------------------------------
\subsection{ปี 2566 ข้อ 26: ยิงประจุเข้าหาประจุอีกตัวจากระยะไกลมาก}

\vspace{1em}

\IfFileExists{img/posn66q26.png}{
\begin{center}
    \includegraphics[width=1\linewidth]{img/posn66q26.png}
\end{center}
}{
\begin{tcolorbox}[title=ตำแหน่งภาพที่แนะนำ]
ให้ตัดภาพข้อสอบปี 2566 ข้อ 26 แล้วบันทึกเป็น
\[
\texttt{img/posn66q26.png}
\]
\end{tcolorbox}
}

\begin{examplebox}{เฉลยละเอียด}
ประจุ $+q$ มวล $m$ ถูกยิงจากระยะไกลมากเข้าหาประจุ $+Q$ มวล $M$ ซึ่งเดิมอยู่นิ่ง ต้องการให้ระยะเข้าใกล้ที่สุดเป็น $D$

ให้พลังงานจลน์เริ่มต้นของประจุ $+q$ เป็น

\[
K=\frac{1}{2}mu^2
\]

ตอนเริ่มต้น พลังงานศักย์ไฟฟ้าเป็นศูนย์เพราะอยู่ไกลมาก

\[
U_i=0
\]

ที่ระยะเข้าใกล้ที่สุด ความเร็วสัมพัทธ์ของสองประจุเป็นศูนย์ คือทั้งสองมีความเร็วเท่ากันเป็นความเร็วของจุดศูนย์กลางมวล

จากโมเมนตัม

\[
mu=(m+M)V
\]

\[
V=\frac{m}{m+M}u
\]

พลังงานจลน์ที่จุดเข้าใกล้ที่สุดคือ

\[
K_f=\frac{1}{2}(m+M)V^2
\]

\[
K_f=\frac{m}{m+M}K
\]

พลังงานศักย์ไฟฟ้าที่ระยะ $D$ คือ

\[
U_f=k_e\frac{qQ}{D}
\]

ใช้อนุรักษ์พลังงาน

\[
K=K_f+U_f
\]

\[
K=\frac{m}{m+M}K+k_e\frac{qQ}{D}
\]

\[
K\left(\frac{M}{m+M}\right)=k_e\frac{qQ}{D}
\]

ดังนั้น

\[
\boxed{K=\frac{m+M}{M}k_e\frac{qQ}{D}}
\]
\end{examplebox}

% ---------------------------------------------------------
\subsection{ปี 2566 ข้อ 27: แรงระหว่างแผ่นตัวเก็บประจุ}

\vspace{1em}

\IfFileExists{img/posn66q27.png}{
\begin{center}
    \includegraphics[width=1\linewidth]{img/posn66q27.png}
\end{center}
}{
\begin{tcolorbox}[title=ตำแหน่งภาพที่แนะนำ]
ให้ตัดภาพข้อสอบปี 2566 ข้อ 27 แล้วบันทึกเป็น
\[
\texttt{img/posn66q27.png}
\]
\end{tcolorbox}
}

\begin{examplebox}{เฉลยละเอียด}
ให้แผ่นตัวเก็บประจุมีประจุ $+Q$ และ $-Q$ พื้นที่แผ่น $A$ และระยะห่าง $d$

ความหนาแน่นประจุผิวคือ

\[
\sigma=\frac{Q}{A}
\]

จากกฎของเกาส์ สนามที่แผ่นประจุหนึ่งแผ่นสร้างมีขนาด

\[
E_{\text{one sheet}}=\frac{\sigma}{2\varepsilon_0}
\]

เมื่อหาแรงบนแผ่นบวก ต้องใช้เฉพาะสนามที่เกิดจากแผ่นลบ เพราะแผ่นบวกไม่ออกแรงต่อตัวเอง

\[
F=QE_{\text{other}}
\]

\[
F=Q\left(\frac{\sigma}{2\varepsilon_0}\right)
\]

แทน $Q=\sigma A$

\[
F=\frac{\sigma^2A}{2\varepsilon_0}
\]

สนามรวมระหว่างแผ่นคือ

\[
E=\frac{\sigma}{\varepsilon_0}
\]

ดังนั้น

\[
F=\frac{1}{2}\varepsilon_0E^2A
\]

ถ้าแหล่งกำเนิดรักษาความต่างศักย์ไว้ที่ $\mathcal{E}$

\[
E=\frac{\mathcal{E}}{d}
\]

จึงได้

\[
F
=
\frac{1}{2}\varepsilon_0A\frac{\mathcal{E}^2}{d^2}
\]

ใช้

\[
C=\varepsilon_0\frac{A}{d}
\]

จะได้

\[
\boxed{F=\frac{C\mathcal{E}^2}{2d}}
\]

ตัวประกอบ $1/2$ เกิดจากการที่แรงบนแผ่นหนึ่งมาจากสนามของอีกแผ่นเท่านั้น ไม่ใช่สนามรวมที่รวมสนามของแผ่นนั้นเอง
\end{examplebox}

% ---------------------------------------------------------
\subsection{ปี 2568 ข้อ 4: สนามไฟฟ้าจากวงแหวนประจุบนแกน}

\vspace{1em}

\IfFileExists{img/posn68q4.png}{
\begin{center}
    \includegraphics[width=1\linewidth]{img/posn68q4.png}
\end{center}
}{
\begin{tcolorbox}[title=ตำแหน่งภาพที่แนะนำ]
ให้ตัดภาพข้อสอบปี 2568 ข้อ 4 แล้วบันทึกเป็น
\[
\texttt{img/posn68q4.png}
\]
\end{tcolorbox}
}

\begin{examplebox}{เฉลยละเอียด}
โจทย์เป็นวงลวดโค้งเป็นวงกลมรัศมี $R$ อยู่ในระนาบ $XOY$ มีประจุรวม $+Q$ จุด $P$ อยู่บนแกน $OZ$ ห่างจากจุด $O$ เป็นระยะ $L$

จากสมมาตร องค์ประกอบสนามในระนาบ $XOY$ หักล้างกันหมด เหลือเฉพาะแกน $z$

สนามบนแกนของวงแหวนคือ

\[
E_z=k_e\frac{QL}{(R^2+L^2)^{3/2}}
\]

ทิศของสนามชี้ออกจากวงแหวนตามแนวจากวงแหวนไปยังจุด $P$ เพราะประจุรวมเป็นบวก

ถ้าจุด $P$ อยู่ที่ $(0,0,-L)$ ทิศของสนามคือแกน $-z$

ดังนั้น

\[
\boxed{\vec{E}=-k_e\frac{QL}{(R^2+L^2)^{3/2}}\hat{k}}
\]
\end{examplebox}


% ---------------------------------------------------------
\subsection{ปี 2568 ข้อ 17: ตัวเก็บประจุหลายตัวต่ออนุกรม}

\vspace{1em}

\IfFileExists{img/posn68q17.png}{
\begin{center}
    \includegraphics[width=1\linewidth]{img/posn68q17.png}
\end{center}
}{
\begin{tcolorbox}[title=ตำแหน่งภาพที่แนะนำ]
ให้ตัดภาพข้อสอบปี 2568 ข้อ 17 แล้วบันทึกเป็น
\[
\texttt{img/posn68q17.png}
\]
\end{tcolorbox}
}

\begin{examplebox}{เฉลยละเอียด}
ตัวเก็บประจุทั้งหมดต่อแบบอนุกรม ดังนั้นประจุบนตัวเก็บประจุทุกตัวมีขนาดเท่ากัน ให้ขนาดประจุนั้นเป็น

\[
q
\]

และความต่างศักย์รวมจากแหล่งกำเนิดคือ

\[
\mathcal{E}
\]

สำหรับตัวเก็บประจุตัวที่ $i$ มี

\[
C_i=i^2C_1
\]

ความจุสมมูลของตัวเก็บประจุอนุกรมหาได้จาก

\[
\frac{1}{C_{\text{eq}}}
=\sum_{i=1}^{n}\frac{1}{C_i}
\]

แทน $C_i=i^2C_1$

\[
\frac{1}{C_{\text{eq}}}
=\sum_{i=1}^{n}\frac{1}{i^2C_1}
\]

\[
\frac{1}{C_{\text{eq}}}
=\frac{1}{C_1}\sum_{i=1}^{n}\frac{1}{i^2}
\]

เมื่อ $n$ มีค่ามาก โจทย์ให้ว่า

\[
1+\frac{1}{2^2}+\frac{1}{3^2}+\frac{1}{4^2}+\cdots+\frac{1}{n^2}\approx\frac{\pi^2}{6}
\]

ดังนั้น

\[
\frac{1}{C_{\text{eq}}}
\approx\frac{1}{C_1}\left(\frac{\pi^2}{6}\right)
\]

จึงได้

\[
C_{\text{eq}}\approx\frac{6C_1}{\pi^2}
\]

ประจุในวงจรอนุกรมคือ

\[
q=C_{\text{eq}}\mathcal{E}
\]

ดังนั้น

\[
q\approx\frac{6C_1\mathcal{E}}{\pi^2}
\]

เพราะตัวเก็บประจุต่ออนุกรมกัน ประจุบนตัวเก็บประจุทุกตัวจึงมีขนาดเท่านี้ รวมถึงตัว $C_i$

\[
\boxed{q_i=\frac{6C_1\mathcal{E}}{\pi^2}}
\]
\end{examplebox}

% ---------------------------------------------------------
\subsection{ปี 2568 ข้อ 19: แรงดึงดูดระหว่างแผ่นตัวเก็บประจุ}

\vspace{1em}

\IfFileExists{img/posn68q19.png}{
\begin{center}
    \includegraphics[width=1\linewidth]{img/posn68q19.png}
\end{center}
}{
\begin{tcolorbox}[title=ตำแหน่งภาพที่แนะนำ]
ให้ตัดภาพข้อสอบปี 2568 ข้อ 19 แล้วบันทึกเป็น
\[
\texttt{img/posn68q19.png}
\]
\end{tcolorbox}
}

\begin{examplebox}{เฉลยละเอียด}
ตัวเก็บประจุแผ่นขนานมีพื้นที่แต่ละแผ่น $A$ อยู่ห่างกัน $d$ และต่อกับแหล่งกำเนิดความต่างศักย์ $V$

เริ่มจากสนามของแผ่นประจุหนึ่งแผ่น โดยใช้พื้นผิวเกาส์รูปกล่องยาคร่อมแผ่น

\[
2E_{\text{sheet}}A
=
\frac{\sigma A}{\varepsilon_0}
\]

ดังนั้น

\[
E_{\text{sheet}}=\frac{\sigma}{2\varepsilon_0}
\]

แรงบนแผ่นหนึ่งเกิดจากสนามของอีกแผ่นเท่านั้น

\[
F=QE_{\text{sheet}}
\]

โดย

\[
Q=\sigma A
\]

จึงได้

\[
F
=(\sigma A)
\left(\frac{\sigma}{2\varepsilon_0}\right)
=
\frac{\sigma^2A}{2\varepsilon_0}
\]

สนามรวมระหว่างแผ่นเป็นผลรวมของสนามจากทั้งสองแผ่น

\[
E=\frac{\sigma}{\varepsilon_0}
\]

ดังนั้น

\[
F=\frac{1}{2}\varepsilon_0E^2A
\]

และเพราะสนามสม่ำเสมอ

\[
E=\frac{V}{d}
\]

จึงได้

\[
\boxed{F=\frac{1}{2}\varepsilon_0A\frac{V^2}{d^2}}
\]

ใช้

\[
C=\varepsilon_0\frac{A}{d}
\]

เขียนได้อีกแบบเป็น

\[
\boxed{F=\frac{CV^2}{2d}}
\]
\end{examplebox}

% ---------------------------------------------------------
\subsection{ปี 2564 ข้อ 8: แรงบนประจุจากวงลวดมีประจุ}

\vspace{1em}

\IfFileExists{img/posn64q8.png}{
\begin{center}
    \includegraphics[width=1\linewidth]{img/posn64q8.png}
\end{center}
}{
\begin{tcolorbox}[title=ตำแหน่งภาพที่แนะนำ]
ให้ตัดภาพข้อสอบปี 2564 ข้อ 8 แล้วบันทึกเป็น
\[
\texttt{img/posn64q8.png}
\]
\end{tcolorbox}
}

\begin{examplebox}{เฉลยละเอียด}
วงลวดรัศมี $R$ มีประจุรวม $q$ กระจายสม่ำเสมอ จุดที่สนใจอยู่บนแกนของวงลวด ห่างจากศูนย์กลางเป็นระยะ $x$

จากสมมาตร องค์ประกอบของสนามไฟฟ้าที่ตั้งฉากกับแกนหักล้างกันหมด เหลือเฉพาะองค์ประกอบตามแกน

สนามไฟฟ้าบนแกนของวงแหวนคือ

\[
E_x=k_e\frac{qx}{(x^2+R^2)^{3/2}}
\]

แรงที่กระทำต่อประจุ $Q$ มีขนาด

\[
F=QE_x
\]

ดังนั้น

\[
F=k_e\frac{qQx}{(x^2+R^2)^{3/2}}
\]

แทน

\[
k_e=\frac{1}{4\pi\varepsilon_0}
\]

จึงได้

\[
\boxed{F=\frac{qQ}{4\pi\varepsilon_0}\frac{x}{(x^2+R^2)^{3/2}}}
\]

\[
\boxed{\text{ตอบ ค}}
\]
\end{examplebox}

% ---------------------------------------------------------
\subsection{ปี 2564 ข้อ 9: ศักย์ไฟฟ้าจากเส้นลวดครึ่งวงกลม}

\vspace{1em}

\IfFileExists{img/posn64q9.png}{
\begin{center}
    \includegraphics[width=1\linewidth]{img/posn64q9.png}
\end{center}
}{
\begin{tcolorbox}[title=ตำแหน่งภาพที่แนะนำ]
ให้ตัดภาพข้อสอบปี 2564 ข้อ 9 แล้วบันทึกเป็น
\[
\texttt{img/posn64q9.png}
\]
\end{tcolorbox}
}

\begin{examplebox}{เฉลยละเอียด}
เส้นลวดครึ่งวงกลมมีรัศมี $R$ และมีประจุรวม $q$ จุด $P$ อยู่บนแกนที่ผ่านศูนย์กลาง ห่างจากศูนย์กลางเป็นระยะ $x$

ทุกชิ้นประจุเล็ก ๆ บนเส้นลวดอยู่ห่างจากจุด $P$ เท่ากัน คือ

\[
r=\sqrt{x^2+R^2}
\]

ศักย์ไฟฟ้าเป็นสเกลาร์ จึงรวมได้โดยตรง

\[
V=k_e\int \frac{dq}{r}
\]

เพราะ $r$ มีค่าเท่ากันสำหรับทุกชิ้นประจุ

\[
V=k_e\frac{1}{\sqrt{x^2+R^2}}\int dq
\]

และ

\[
\int dq=q
\]

ดังนั้น

\[
V=k_e\frac{q}{\sqrt{x^2+R^2}}
\]

หรือ

\[
\boxed{V=\frac{q}{4\pi\varepsilon_0\sqrt{x^2+R^2}}}
\]

\[
\boxed{\text{ตอบ ก}}
\]
\end{examplebox}

% ---------------------------------------------------------
\subsection{ปี 2564 ข้อ 10: พลังงานความร้อนหลังต่อตัวเก็บประจุเข้าด้วยกัน}

\vspace{1em}

\IfFileExists{img/posn64q10.png}{
\begin{center}
    \includegraphics[width=1\linewidth]{img/posn64q10.png}
\end{center}
}{
\begin{tcolorbox}[title=ตำแหน่งภาพที่แนะนำ]
ให้ตัดภาพข้อสอบปี 2564 ข้อ 10 แล้วบันทึกเป็น
\[
\texttt{img/posn64q10.png}
\]
\end{tcolorbox}
}

\begin{examplebox}{เฉลยละเอียด}
เดิมตัวเก็บประจุ $C_1$ มีประจุ $q_0$ ส่วน $C_2$ ยังไม่มีประจุ เมื่อปิดสวิตช์ ประจุจะไหลจนตัวเก็บประจุทั้งสองมีความต่างศักย์เท่ากัน

ประจุรวมคงที่

\[
q_{\text{total}}=q_0
\]

ความจุรวมหลังต่อร่วมกันคือ

\[
C_{\text{total}}=C_1+C_2
\]

ดังนั้นความต่างศักย์สุดท้ายคือ

\[
V_f=\frac{q_0}{C_1+C_2}
\]

พลังงานเริ่มต้นอยู่ใน $C_1$ เท่านั้น

\[
U_i=\frac{q_0^2}{2C_1}
\]

พลังงานสุดท้ายในตัวเก็บประจุทั้งสองคือ

\[
U_f=\frac{1}{2}(C_1+C_2)V_f^2
\]

แทน $V_f$

\[
U_f=\frac{1}{2}(C_1+C_2)\left(\frac{q_0}{C_1+C_2}\right)^2
\]

\[
U_f=\frac{q_0^2}{2(C_1+C_2)}
\]

พลังงานที่หายไปกลายเป็นความร้อนในตัวต้านทาน

\[
Q_{\text{heat}}=U_i-U_f
\]

\[
Q_{\text{heat}}=\frac{q_0^2}{2C_1}-\frac{q_0^2}{2(C_1+C_2)}
\]

\[
Q_{\text{heat}}
=\frac{1}{2}\left(\frac{C_2}{C_1(C_1+C_2)}\right)q_0^2
\]

หรือ

\[
\boxed{Q_{\text{heat}}
=\frac{1}{2}\left(\frac{C_2/C_1}{C_1+C_2}\right)q_0^2}
\]

\[
\boxed{\text{ตอบ ค}}
\]
\end{examplebox}

% ---------------------------------------------------------
\subsection{ปี 2564 ข้อ 16: พลังงานศักย์ไฟฟ้าของระบบประจุบนวงกลม}

\vspace{1em}

\IfFileExists{img/posn64q16.png}{
\begin{center}
    \includegraphics[width=1\linewidth]{img/posn64q16.png}
\end{center}
}{
\begin{tcolorbox}[title=ตำแหน่งภาพที่แนะนำ]
ให้ตัดภาพข้อสอบปี 2564 ข้อ 16 แล้วบันทึกเป็น
\[
\texttt{img/posn64q16.png}
\]
\end{tcolorbox}
}

\begin{examplebox}{เฉลยละเอียด}
ระบบมีประจุ $+2e$ ที่ศูนย์กลาง และมีประจุ $-e$ สองตัวอยู่บนวงกลมรัศมี $R$ โดยทำมุมกัน $\theta$

พลังงานศักย์ไฟฟ้ารวมของระบบต้องรวมทุกคู่

\[
U=\sum_{i<j}k_e\frac{q_iq_j}{r_{ij}}
\]

คู่ระหว่าง $+2e$ กับ $-e$ แต่ละตัวอยู่ห่างกัน $R$ จึงให้พลังงาน

\[
U_1=k_e\frac{(+2e)(-e)}{R}=-\frac{2k_ee^2}{R}
\]

มีสองคู่แบบนี้ จึงรวมเป็น

\[
U_{\text{center}}=-\frac{4k_ee^2}{R}
\]

ต่อไปพิจารณาคู่ระหว่างประจุ $-e$ ทั้งสองตัว ระยะระหว่างกันคือคอร์ดของวงกลม

\[
d=2R\sin\frac{\theta}{2}
\]

ดังนั้น

\[
U_{--}=k_e\frac{(-e)(-e)}{d}
=\frac{k_ee^2}{2R\sin(\theta/2)}
\]

พลังงานรวมคือ

\[
U=-\frac{4k_ee^2}{R}+\frac{k_ee^2}{2R\sin(\theta/2)}
\]

พจน์แรกเป็นค่าคงที่ ส่วนพจน์ที่สองจะน้อยที่สุดเมื่อ

\[
\sin\frac{\theta}{2}
\]

มีค่ามากที่สุด ซึ่งเกิดเมื่อ

\[
\frac{\theta}{2}=90^\circ
\]

ดังนั้น

\[
\theta=180^\circ
\]

\[
\boxed{\text{ตอบ ง}}
\]
\end{examplebox}

% ---------------------------------------------------------
\subsection{ปี 2565 ข้อ 3: ทิศของสนามไฟฟ้าลัพธ์จากสองประจุ}

\vspace{1em}

\IfFileExists{img/posn65q3.png}{
\begin{center}
    \includegraphics[width=1\linewidth]{img/posn65q3.png}
\end{center}
}{
\begin{tcolorbox}[title=ตำแหน่งภาพที่แนะนำ]
ให้ตัดภาพข้อสอบปี 2565 ข้อ 3 แล้วบันทึกเป็น
\[
\texttt{img/posn65q3.png}
\]
\end{tcolorbox}
}

\begin{examplebox}{เฉลยละเอียด}
ต้องหาทิศของสนามไฟฟ้าลัพธ์ที่จุด $(a,a)$ จากประจุ $+2Q$ ที่ $(0,0)$ และประจุ $-Q$ ที่ $(2a,0)$

ระยะจากประจุทั้งสองถึงจุด $(a,a)$ เท่ากัน คือ

\[
r=a\sqrt{2}
\]

สนามจากประจุ $+2Q$ มีขนาด

\[
E_1=k_e\frac{2Q}{(a\sqrt{2})^2}
\]

\[
E_1=k_e\frac{Q}{a^2}
\]

ทิศของ $\vec{E}_1$ ชี้จากจุดกำเนิดไปยัง $(a,a)$ จึงมีองค์ประกอบ

\[
E_{1x}=\frac{E_1}{\sqrt{2}},\qquad E_{1y}=\frac{E_1}{\sqrt{2}}
\]

สนามจากประจุ $-Q$ มีขนาด

\[
E_2=k_e\frac{Q}{(a\sqrt{2})^2}
\]

\[
E_2=k_e\frac{Q}{2a^2}=\frac{E_1}{2}
\]

เพราะเป็นประจุลบ สนามที่จุด $(a,a)$ ชี้เข้าหาประจุ $-Q$ คือชี้ลงทางขวา

\[
E_{2x}=\frac{E_2}{\sqrt{2}},\qquad E_{2y}=-\frac{E_2}{\sqrt{2}}
\]

รวมองค์ประกอบ

\[
E_x=\frac{E_1}{\sqrt{2}}+\frac{E_2}{\sqrt{2}}
=\frac{3E_1}{2\sqrt{2}}
\]

\[
E_y=\frac{E_1}{\sqrt{2}}-\frac{E_2}{\sqrt{2}}
=\frac{E_1}{2\sqrt{2}}
\]

ดังนั้น

\[
\tan\theta=\frac{E_y}{E_x}
\]

\[
\tan\theta=\frac{1}{3}
\]

จึงได้

\[
\boxed{\theta=\tan^{-1}\left(\frac{1}{3}\right)}
\]

\[
\boxed{\text{ตอบ ง}}
\]
\end{examplebox}

% ---------------------------------------------------------
\subsection{ปี 2565 ข้อ 4: ทรงกลมนำไฟฟ้ากับเปลือกทรงกลมต่อลงดิน}

\vspace{1em}

\IfFileExists{img/posn65q4.png}{
\begin{center}
    \includegraphics[width=1\linewidth]{img/posn65q4.png}
\end{center}
}{
\begin{tcolorbox}[title=ตำแหน่งภาพที่แนะนำ]
ให้ตัดภาพข้อสอบปี 2565 ข้อ 4 แล้วบันทึกเป็น
\[
\texttt{img/posn65q4.png}
\]
\end{tcolorbox}
}

\begin{examplebox}{เฉลยละเอียด}
ทรงกลมโลหะด้านในมีรัศมี $R_1$ และมีประจุ $-Q$ ส่วนทรงกลมโลหะด้านนอกมีรัศมี $R_2$ และต่อลงดิน

ต้องการหาว่าศักย์ไฟฟ้าของทรงกลมโลหะด้านนอกสูงกว่าด้านในเท่าไร นั่นคือ

\[
V_{\text{out}}-V_{\text{in}}
\]

ในบริเวณระหว่างทรงกลมทั้งสอง สนามไฟฟ้ามีขนาดเหมือนสนามจากประจุ $-Q$ ที่ศูนย์กลาง

\[
E_r=-k_e\frac{Q}{r^2}
\]

ใช้ความสัมพันธ์

\[
V_{\text{out}}-V_{\text{in}}
=-\int_{R_1}^{R_2} \vec{E}\cdot d\vec{r}
\]

แทนค่า

\[
V_{\text{out}}-V_{\text{in}}
=-\int_{R_1}^{R_2}\left(-k_e\frac{Q}{r^2}\right)dr
\]

\[
V_{\text{out}}-V_{\text{in}}
=k_eQ\int_{R_1}^{R_2}\frac{dr}{r^2}
\]

\[
V_{\text{out}}-V_{\text{in}}
=k_eQ\left(\frac{1}{R_1}-\frac{1}{R_2}\right)
\]

ดังนั้น

\[
\boxed{V_{\text{out}}-V_{\text{in}}
=\frac{Q}{4\pi\varepsilon_0}\left(\frac{1}{R_1}-\frac{1}{R_2}\right)}
\]

\[
\boxed{\text{ตอบ ข}}
\]
\end{examplebox}

% ---------------------------------------------------------
\subsection{ปี 2565 ข้อ 7: ประจุสองตัวโคจรรอบกันด้วยแรงคูลอมบ์}

\vspace{1em}

\IfFileExists{img/posn65q7.png}{
\begin{center}
    \includegraphics[width=1\linewidth]{img/posn65q7.png}
\end{center}
}{
\begin{tcolorbox}[title=ตำแหน่งภาพที่แนะนำ]
ให้ตัดภาพข้อสอบปี 2565 ข้อ 7 แล้วบันทึกเป็น
\[
\texttt{img/posn65q7.png}
\]
\end{tcolorbox}
}

\begin{examplebox}{เฉลยละเอียด}
ประจุ $+Q$ และ $-Q$ มีมวลเท่ากัน $m$ และอยู่ห่างกันคงที่ $D$ โดยโคจรรอบกันด้วยอัตราเร็ว $v$

เพราะมวลเท่ากัน จุดศูนย์กลางการโคจรอยู่กึ่งกลางระหว่างสองประจุ ดังนั้นรัศมีการเคลื่อนที่ของแต่ละประจุคือ

\[
r=\frac{D}{2}
\]

แรงคูลอมบ์เป็นแรงสู่ศูนย์กลาง

\[
F_e=k_e\frac{Q^2}{D^2}
\]

และ

\[
F_c=\frac{mv^2}{D/2}
\]

ตั้งให้เท่ากัน

\[
k_e\frac{Q^2}{D^2}=\frac{2mv^2}{D}
\]

จัดรูป

\[
\frac{mv^2D}{Q^2}=\frac{k_e}{2}
\]

แทน

\[
k_e=\frac{1}{4\pi\varepsilon_0}
\]

จึงได้

\[
\boxed{\frac{mv^2D}{Q^2}=\frac{1}{8\pi\varepsilon_0}}
\]

\[
\boxed{\text{ตอบ ข}}
\]
\end{examplebox}

% ---------------------------------------------------------
\subsection{ปี 2565 ข้อ 17: ความจุของตัวเก็บประจุทรงกระบอกซ้อนกัน}

\vspace{1em}

\IfFileExists{img/posn65q17.png}{
\begin{center}
    \includegraphics[width=1\linewidth]{img/posn65q17.png}
\end{center}
}{
\begin{tcolorbox}[title=ตำแหน่งภาพที่แนะนำ]
ให้ตัดภาพข้อสอบปี 2565 ข้อ 17 แล้วบันทึกเป็น
\[
\texttt{img/posn65q17.png}
\]
\end{tcolorbox}
}

\begin{examplebox}{เฉลยละเอียด}
รูป ก เป็นตัวเก็บประจุแผ่นขนาน ซึ่งมีสูตร

\[
C=\varepsilon_0\frac{A}{d}
\]

รูป ข เป็นตัวเก็บประจุทรงกระบอกซ้อนกัน มีรัศมีด้านในประมาณ $r$ ความยาว $l$ และระยะห่างระหว่างผิวตัวนำเท่ากับ $d$

ถ้า $d$ เล็กเมื่อเทียบกับ $r$ เรามองผิวทรงกระบอกเหมือนแผ่นขนานที่ถูกม้วนเป็นทรงกระบอกได้

พื้นที่ผิวที่เผชิญกันประมาณเป็นพื้นที่ผิวด้านข้างของทรงกระบอกด้านใน

\[
A\approx 2\pi r l
\]

แทนในสูตรแผ่นขนาน


\[
C\approx \varepsilon_0\frac{2\pi rl}{d}
\]

ดังนั้น

\[
\boxed{C\approx \frac{2\pi\varepsilon_0rl}{d}}
\]

\[
\boxed{\text{ตอบ ข}}
\]
\end{examplebox}